% \documentclass[acmsmall]{acmart}
\documentclass[manuscript, screen, review]{acmart}


%% \BibTeX command to typeset BibTeX logo in the docs
\AtBeginDocument{%
  \providecommand\BibTeX{{%
    Bib\TeX}}}



%% Rights management information. This information is sent to you
%% when you complete the rights form. These commands have SAMPLE
%% values in them; it is your responsibility as an author to replace
%% the commands and values with those provided to you when you
%% complete the rights form.
\setcopyright{acmlicensed}
\copyrightyear{2026}
\acmYear{2026}
\acmDOI{XXXXXXX.XXXXXXX}

%% These commands are for a JOURNAL article.
\acmJournal{FAC}
\acmVolume{0}
\acmNumber{0}
\acmArticle{1}
\acmMonth{1}

%% Submission ID.
%% Use this when submitting an article to a sponsored event. You'll
%% receive a unique submission ID from the organizers of the event,
%% and this ID should be used as the parameter to this command.
%%\acmSubmissionID{123-A56-BU3}

%% For managing citations, it is recommended to use bibliography
%% files in BibTeX format.
%%
%% You can then either use BibTeX with the ACM-Reference-Format style,
%% or BibLaTeX with the acmnumeric or acmauthoryear styles, that include
%% support for advanced citation of software artefact from the
%% biblatex-software package, also separately available on CTAN.
%%
%% Look at the sample-*-biblatex.tex files for templates showcasing
%% the biblatex styles.

%% The majority of ACM publications use numbered citations and
%% references. The command \citestyle{authoryear} switches to the
%% "author year" style.
%%
%% If you are preparing content for an event
%% sponsored by ACM SIGGRAPH, you must use the "author year" style of
%% citations and references.
%% Uncommenting the next command will enable that style.
%%\citestyle{acmauthoryear}

%%
%% end of the preamble, start of the body of the document source.

\usepackage{enumerate}

\input NEW-commands-FAC.tex

\newcommand{\AUX}{\mathit{AUX}}

\newtheorem{note}{Note}  % was not present in ACM style


\begin{document}

%%
%% The "title" command has an optional parameter,
%% allowing the author to define a "short title" to be used
%% in page headers.

\title[Termination of Loops]%
 {Proof Rules for Termination of Loops: A Comparative Analysis}

%%
%% The "author" command and its associated commands are used to define
%% the authors and their affiliations.
%% Of note is the shared affiliation of the first two authors, and the
%% "authornote" and "authornotemark" commands
%% used to denote shared contribution to the research.
\author{Krzysztof R. Apt}
\affiliation{%
  \institution{Centrum voor Wiskunde en Informatica}
  \city{Amsterdam}
  \country{The Netherlands,}
\institution{MIMUW}  % How to do second affiliation ? 
\city{Warsaw}
\country{Poland}
}
\email{apt@cwi.nl}
\orcid{0000-0002-1332-4229}

\author{Frank S. de Boer}
\affiliation{%
  \institution{Centrum voor Wiskunde en Informatica}
  \city{Amsterdam}
  \country{The Netherlands}
}
\email{F.S.de.Boer@cwi.nl}
\orcid{???}

\author{Ernst-R\"{u}diger Olderog}
\affiliation{%
  \institution{Carl von Ossietzky Universit\"{a}t Oldenburg}
  \city{Oldenburg}
  \country{Germany}
}
\email{olderog@informatik.uni-oldenburg.de}
\orcid{0009-0009-6314-1269}



%%
%% By default, the full list of authors will be used in the page
%% headers. Often, this list is too long, and will overlap
%% other information printed in the page headers. This command allows
%% the author to define a more concise list
%% of authors' names for this purpose.
\renewcommand{\shortauthors}{K.R. Apt, F.S. de Boer and E.-R. Olderog}


\begin{abstract}
  We provide a comparative analysis of different proof rules for
  proving termination of \textbf{while} programs and show their
  proof-theoretic equivalence.  This involves a proof-theoretic
  analysis of various auxiliary proof rules in Hoare's logic.  By
  discussing representations of proofs in the form of proof outlines,
  we reveal differences between these equivalent proof rules when used
  in practice.  We also address applications in the context of the
  paradigm of design by contract.
\end{abstract}

%%
%% The code below is generated by the tool at http://dl.acm.org/ccs.cfm.
%% Please copy and paste the code instead of the example below.
%%
\begin{CCSXML}
<ccs2012>
<concept>
<concept_id>10003752.10003790.10011741</concept_id>
<concept_desc>Theory of computation~Hoare logic</concept_desc>
<concept_significance>500</concept_significance>
</concept>
</ccs2012>
\end{CCSXML}

\ccsdesc[500]{Theory of computation~Hoare logic}

%% Keywords. The author(s) should pick words that accurately describe
%% the work being presented. Separate the keywords with commas.
\keywords{Program verification, Termination, Analysis of proof rules}

\received{Day Month 2026}
% \received[revised]{?? ???? 2025}
% \received[accepted]{?? ???? 2025}

% \begin{teaserfigure}
%   \includegraphics[width=\textwidth]{???}
%   \caption{figure caption}
%   \Description{figure description}
% \end{teaserfigure}

%%
%% This command processes the author and affiliation and title
%% information and builds the first part of the formatted document.
\maketitle




\section{Introduction}

In his seminal paper \cite{Hoa69}, Hoare introduced an axiomatic method of reasoning
about correctness of \textbf{while} programs, now called Hoare's logic. 
It is based on correctness formulas
$\HT{p}{S}{q}$, where $S$ is a program and $p$ and $q$ are assertions
(logical formulas), with the interpretation
   \begin{quote}
   ``If the assertion $p$ is true before initiation of a program $S$,
     then the assertion $q$ will be true on its completion.''
   \end{quote}
In this context $p$ is referred to as a \emph{precondition} and $q$ as 
a \emph{postcondition} of $S$.

However, in contrast to Floyd's earlier paper \cite{Flo67a} that dealt
with correctness of flowchart programs, program termination was not
addressed. To stress this difference one distinguishes now between
\emph{partial correctness} that only focuses on the delivery of
correct results, and \emph{total correctness}, that in addition
stipulates that the program terminates.  So the original proposal of
Hoare dealt with partial correctness.

All approaches to proving program termination within Hoare’s logic formalize 
Floyd’s \cite{Flo67a} observation that
    \begin{quote}
    ``Proofs of termination are dealt with by showing that each step of a program 
      decreases some entity which cannot decrease indefinitely.”
    \end{quote}

The first extension of Hoare's logic
to deal with total correctness is due to~\cite{MP74}. 
Since then substantially simpler proof rules were proposed.
In these proof rules variables that range over natural numbers are used.
An appropriate relative completeness result, see for example \cite{ABO09}, 
shows that variables ranging over more general well-founded orderings 
are not needed.

Termination continues to be a relevant and vibrant topic in program
analysis, see for example~\cite{CookPR11} and the Annual International
Termination Competition\footnote{
\texttt{https://termination-portal.org/wiki/Termination\_Competition}}.
The latter comprises various competition categories, 
for instance proving termination of 
C programs, Java bytecode programs, logic programs,
functional programs, and term rewriting systems.
Here we focus on the shape of termination proofs in the context of
Hoare's logic. 
To this end, we investigate three natural proof rules 
from the proof-theoretic point of view.
More specifically, we study {\bf while} programs with Hoare's 
original proof system for program correctness, in which the well-known
proof rule for partial correctness of the {\bf while} loops, which we
call the LOOP I rule, is replaced by a suitable proof rule for
establishing total correctness.  We analyze three versions of such a
proof rule. 


\medskip


\begin{itemize}

\item 
% The LOOP III rule allows us to document proofs as \emph{proof
%   outlines}, which is in particular useful when arguing about the
%   interference freedom of component proofs in the context of parallel
%   programs.
The LOOP II rule involves a single, monolithic proof which
  establishes an invariant, termination of the body of the \textbf{while}
  loop, and which ensures that the body itself cannot be executed
  infinitely often by a so-called bound function.
% \item The LOOP IV rule is particularly well-suited when dealing with nested loops
%       because it modularizes the correctness proof of the outer loop from
%       the ones of the inner loops. It is a \emph{hybrid} proof rule with
%       premises referring to proof systems for both partial and total correctness.

\item The LOOP III rule achieves a separation of the reasoning about
  the invariant and the bound function.

\item
The LOOP IV rule not only achieves a separation of the reasoning about
the invariant and the bound function, it additionally involves
a separate proof of the termination the body itself.
As such it is particularly well-suited when dealing with nested loops
because it modularizes the correctness proof of the outer loop from
the ones of the inner loops. It is a \emph{hybrid} proof rule with
premises referring to proof systems for both partial and total correctness.

% \item The LOOP V rule, like the LOOP II rule, is a monolithic rule
% which formalizes reasoning  about termination of a \textbf{while} loop in a single proof.
% However, instead  of an explicit representation of the bound function
% it uses  a parameter of the invariant itself.

\end{itemize}

Depending on the choice of the loop rule, we obtain proof systems that
we refer to by II, III, and IV, respectively.  We show that these
three proof systems are equivalent in the sense that every proof of a
correctness formula $\HT{p}{S}{q}$ carried out in one of these systems
can be effectively transformed into a correctness proof in any other
of these proof systems.  This result is obtained by a detailed
proof-theoretic analysis of the four loop rules in the context of
these proof systems. In this analysis we make use of
auxiliary proof rules which we show to be admissible in the proof systems.

Even though Hoare's logic has been extensively studied (see for
example our survey \cite{AO19}), little work has been done on the
analysis of proofs in Hoare's logic. We are familiar with only three
references, \cite{Apt81a}, \cite{GR16}, and \cite{Tiu02}, in which
transformations of proofs in a Hoare logic are discussed.

While these proof rules are equivalent, their use and representation
in the form of proof outlines, which are programs annotated by
assertions, differs. We illustrate this by analyzing a termination
proof of a program with nested loops. We also address applications in
the context of assertions used as annotations in the design by
contract paradigm.
Finally, we show how to extend the equivalence result to a larger class of programs
that access a library of procedures.


\section{LOOP rules}

We shall be concerned here with the language, the formulas of
which are either first-order formulas, called \bfe{assertions}, or
\bfe{correctness formulas}, which are constructs of the form $\HT{p}{S}{q}$,
where $p$ and $q$ are assertions and $S$ is a \textbf{while} program.

We assume two disjoint sets of variables: a set \emph{lvar} of
\bfe{logical} variables and a set \emph{pvar} of \bfe{program}
variables. Programs refer only to program variables, while assertions
may refer to both types of variables.  Logical variables are the only
variables that can be quantified over. They are used in assertions as
\bfe{freeze} variables.
To keep things simple we assume that all program and logical variables
are simple variables of  type \textbf{integer}.

Below we denote by $\mathit{free}(p)$
the set of free variables of the assertion $p$ and by $var(t)$, $var(B)$, and $var(S)$ 
the set of variables that appear in the
expression $t$, the Boolean expression $B$, and the program $S$, respectively.
%
% ADDED NOTATION SUBSTITUTION
% 
% \NI KRA (28-01): I WOULD LIKE TO SUGGEST TO DELETE
% THE SECOND SENTENCE: IT IS A STANDARD DEFINITION. ALSO WE NEVER HAVE
% TO PERFORM SUCH A RENAMING. 
%
We denote by $p[x:=t]$ the result of substituting all free occurrences
of the variable $x$ in $p$ by the expression $t$. 

% The standard
% definition takes care by means of renaming of bound variables that
% after the replacement no variable of $t$ becomes bound in $p$.


We shall consider in total four proof systems concerned with the
correctness formulas. They only differ in the used LOOP rule.

Proof system I denotes the customary proof system allowing us to prove 
partial correctness of {\bf while} programs.  
Its axioms and proof rules, due to \cite{Hoa69}, are listed in Appendix~\ref{sec:I}.
Its LOOP rule has the following form:
\III

\NI
RULE LOOP I
\[ 
 \frac{ \HT{p \A B}{S}{p}                }
      { \HT{p}{\WDD{B}{S}}{p \A \neg B}  }
\]
\III


In the proof system II, this rule  is replaced by
\III

\NI
RULE LOOP II
\[
 \begin{array}{l}
  \HT{p \A B \A t=z}{S}{p \A t<z},          \\
  p \ra t \geq 0                           \\
  [-\medskipamount]
  \hrulefill                                \\
  \HT{p}{\WDD{B}{S}}{p \A \neg\ B}
 \end{array}
\]
where $t$ is an integer expression such that  $var(t)\subseteq \mbox{\it pvar}$
and $z$ is a \emph{logical} integer variable that does not appear in $p$.

% KRA: (28-01) I WOULD LIKE TO SUGGEST TO DELETE 'free' (IT WAS '(free)'). NOTHING IS LOST
% BY SUCH A STRONGER REQUIREMENT.  

\III

In the proof system III the  LOOP I rule is replaced by
\III

\NI
RULE \label{rul:loop2} LOOP III
\[
 \begin{array}{l}
  \HT{p \A B}{S}{p},                      \\
  \HT{p \A B \A t=z}{S}{t<z},             \\
  p \ra t \geq 0                         \\
  [-\medskipamount]
  \hrulefill                              \\
  \HT{p}{\WDD{B}{S}}{p \A \neg\ B}
 \end{array}
\]
%
\NI
where $t$ and $z$ are as above.

\III

In the context of the LOOP rules discussed here, 
the assertion $p$ is called the \bfe{loop invariant} and the expression $t$
is called the \bfe{bound function}. It provides an estimate how
many iterations the loop will still perform before termination.

We shall consider next the following hybrid rule that combines provability in two
proof systems.
\III

\NI
RULE LOOP IV
\[
 \begin{array}{l}
  \vdash_{\rm{I}}  \HT{p \A B}{S}{p},             \\
  \vdash_{\rm{I}}   \HT{p \A B \A t=z}{S}{t<z},   \\
  \HT{p \A B}{S}{\mathbf{true}},                  \\
  p \ra t \geq 0                                 \\
  [-\medskipamount]
  \hrulefill                                      \\
  \HT{p}{\WDD{B}{S}}{p \A \neg\ B}
 \end{array}
\]
%
\NI
where $t$ and $z$ are as above.
\III

Proof system~IV is obtained from proof system~I by replacing the LOOP
I rule by the LOOP~IV rule. The use of two forms of provability in the
premises of this rule can be circumvented by the following
modification of the notation.  Denote the correctness
formulas in the sense of partial correctness by $\HT{p}{S}{q}$ and in
the sense of total correctness by $\HTT{p}{S}{q}$. Then combine the
proof system~I with the proof system in which the axioms and proof
rules of I except the LOOP~I rule are rewritten using the
$\HTT{p}{S}{q}$ syntax. Finally, add to this proof system the LOOP IV
rewritten as follows:
\III

\[
 \begin{array}{l}
  \HT{p \A B}{S}{p},                      \\
  \HT{p \A B \A t=z}{S}{t<z},             \\
  \HTT{p \A B}{S}{\mathbf{true}},         \\
  p \ra t \geq 0                         \\
  [-\medskipamount]
  \hrulefill                              \\
  \HTT{p}{\WDD{B}{S}}{p \A \neg\ B}
 \end{array}
\]

\NI
where $t$ and $z$ are as above.
\III

In what follows we use the original formulation of this rule, 
as it will not lead to any ambiguities.

% The final LOOP V rule introduces an additional parameter
% in the loop invariant instead of a separate bound function.

% \III

% \NI
% RULE LOOP V
% \[
%  \begin{array}{l}
%  \HT{p(x)\wedge B}{S}{\exists y: p(y)\wedge y< x},          \\
%   p(x)\rightarrow x\geq 0                          \\
%   [-\medskipamount]
%   \hrulefill                                \\
%   \HT{\exists x:p(x)}{\WDD{B}{S}}{(\exists x:p(x)) \A \neg\ B}
%  \end{array}
% \]

% \NI
% where $x$ is an integer variable which  does not appear
% in $\WDD{B}{S}$, and $p(y)$ denotes the result of replacing
% $x$ in $p(x)$ by the fresh integer variable $y$.

% Note that the outer brackets in $(\exists x:p(x))$ can be dropped thanks to the assumption about the variable $x$.


The main aim of this paper is to prove the following result.

\begin{theorem}[Equivalence] \label{thm:abo}
The proof systems II-IV are equivalent.  
\end{theorem}

Here, of course, equivalence of the proof systems means that the same
formulas can be proved in each of them.



\section{Discussion}

Before we proceed with the proof of the Equivalence Theorem~\ref{thm:abo} let us
discuss the introduced rules and some of their variants.  It is
difficult to determine where the LOOP II rule was introduced first. We
mentioned it in \cite{ABO90}. It also appears in \cite[page 64]{Rey98}
and in \cite[page 151]{Almeida11}.  The LOOP III rule was introduced
in \cite{OG76a}.  It corresponds to Dijkstra's modification of his
weakest precondition semantics proposed in \cite{EWD:EWD573} and
reproduced as %\cite{Dij82}.
\cite{EWD:EWD573pub}.  We used it in our book \cite{ABO09}.

The LOOP IV rule was introduced in \cite{ABO25}.
It formalizes the following intuition. In order to prove the
termination of a \textbf{while} loop it suffices to find a loop
invariant and a bound function such that

\begin{enumerate}[(i)]
 \item the loop invariant is maintained by each loop body execution, 
       in the sense of partial correctness,
 \item each loop body execution decreases the bound function, 
       also in the sense of partial correctness,
 \item the loop body terminates, and
 \item the loop invariant implies that the bound function remains non-negative.
\end{enumerate}

This rule supports modular reasoning about program correctness
by separating the premises into partial correctness and termination
properties.  This is of particular relevance in the presence of nested
loops, i.e., when the loop body contains inner loops.  Then (i)
establishes only the partial correctness of the loop body, whereas its
termination is relegated to (iii). That the loop body can be iterated
only finitely often is established in (ii) in combination with (iv).
Note that for loop bodies without inner loops, partial and total
correctness coincide. In this case the third premise,
$\HT{p \A B}{S}{\mathbf{true}}$, can then be dropped, and we arrive at
the LOOP III rule. So the hybrid form of this rule arises only for
nested loops.

% Parameterized loop invariants were first used in a rule proposed
% in \cite{Harel79}, where with each loop iteration the parameter was
% supposed to decrement exactly by~1.  The above LOOP V rule allows us
% to generalize proofs of termination to arbitrary well-founded domains.
% Such a generalization is for example required for reasoning about
% termination in presence of \emph{countable non-determinism}
% \cite{AptP86}.  

It is easy to manufacture a number of alternative LOOP rules. 
Note for example that the LOOP II rule results from combining the first two
premises of the LOOP III into one.  An analogous modification can be
carried out in the case of the LOOP IV rule. 

% In \cite{AptP86} a different version of LOOP V rule was used, namely

% \[
%  \begin{array}{l}
%    p(x) \wedge x > 0 \ra B,                         \\
%  \HT{p(x)\wedge x > 0}{S}{\exists y: p(y)\wedge y< x},          \\
%    p(0) \ra \neg B,                         \\ 
%   [-\medskipamount]
%   \hrulefill                                \\
%   \HT{\exists x:p(x)}{\WDD{B}{S}}{p(0)}
%  \end{array}
% \]

% (In this rule, to deal with unbounded nondeterminism, it was stipulated
% that the variable $x$ ranges over ordinals instead over integers.)
% We found our LOOP rule V simpler to reason with.

% In turn, the following version of LOOP V results from separating the
% reasoning about termination along the lines of LOOP IV rule.

% \[
%  \begin{array}{l}
%  \vdash_{\rm{I}}\HT{p(x)\wedge B}{S}{\exists y: p(y)\wedge y< x},          \\
%  \HT{(\exists x:p) \wedge B}{S}{\T},\\
%   p(x)\rightarrow x\geq 0                          \\
%   [-\medskipamount]
%   \hrulefill                                \\
%   \HT{\exists x:p(x)}{\WDD{B}{S}}{(\exists x:p(x)) \A \neg\ B}
%  \end{array}
% \]


KRA (6-03-26):BEGIN REVISED (and moved here)

Further, the last premise of each of the LOOP rules II-IV and their
variants can also be modified, by considering in each case the easier
to establish implication $p \A B \ra t \geq 0$ instead of
$p \ra t \geq 0$.  In fact, in \cite[page 64]{Rey98} and in \cite[page
151]{Almeida11} this variant of the LOOP II is used. We call the
corresponding rules the LOOP $B$ rules.

We shall clarify their proof-theoretic status in Appendix \ref{sec:B}.

KRA: END 



\section{Proving Equivalence}

A direct proof of the Equivalence Theorem~\ref{thm:abo} would require a detailed
analysis of the shape of proofs. For example, in the case of the proof
systems II and III we would have to construct in II a proof of the premise
$ \HT{p \A B \A t=z}{S}{p \A t<z}$ of the LOOP II rule from the proofs
of the premises 
\[
 \HT{p \A B}{S}{p} \mbox{ and } \HT{p \A B \A t=z}{S}{t<z}
\]
of the LOOP III rule, which looks daunting.

Therefore we proceed differently. First we introduce a set of
auxiliary rules using which we can deduce each LOOP rule from the
other and then we show how we get get rid of these auxiliary rules.
To explain this approach we introduce first the relevant notions from
proof theory.


\subsection{Admissibility and Derivability}

Assume a given language that is determined by its set of formulas.  In
what follows we assume that all considered axioms, proof rules, and
proof systems are concerned with the same language.

Given a proof system $T$ and two sequences of formulas
$\phi_1, \LL, \phi_m$ and $\varphi_1, \LL, \varphi_n$ we write
\[
  \phi_1, \LL, \phi_m \vdash_T \varphi_1, \LL, \varphi_n
\]
to denote the fact that each formula $\varphi_i$ can be proved in
$T$ using as additional axioms the formulas
$\phi_1, \LL, \phi_m$.  We also use this notation when the sequence
$\phi_1, \LL, \phi_m$ is empty and when $T$ is a set of proof
rules.

Given a proof system $T$ and a proof rule (in short: a rule) $R$ we
abbreviate $T \cup \Cu{R}$ to $T \cup R$. Given two proof systems $T_1$
and $T_2$ we abbreviate the statement that they are equivalent to
$T_1 \approx T_2$ and the statement that every formula provable in
$T_1$ is also provable in $T_2$ to $T_1 \preceq T_2$.

\begin{definition}(See, e.g.~\cite[page 47]{Pla13}.)
  
  Consider a rule
  \[ 
  R \qquad \frac{\varphi_1,\LL,\varphi_k}{\varphi}
\]
and a proof system $T$. $R$ is called
\begin{itemize}
\item \emph{admissible in $T$} if 
\[
 \vdash_T \varphi_1, \LL, \vdash_T \varphi_k
 \mbox{ implies }
 \vdash_T \varphi,
\]

\item \emph{derivable in $T$} if 
\[
 \varphi_1, \LL, \varphi_k \vdash_T \varphi.
\]

\end{itemize}

\end{definition}

So for an admissible rule a proof in $T$ of its conclusion 
that is manufactured by \emph{manipulating} $k$ proofs in $T$ 
of its premises can be condensed into \emph{one step}. 
In turn, a derivable rule allows us to condense a proof in $T$ of
its conclusion from its premises into \emph{one step}.

Both types of rules do not increase the power of the proof system $T$
but serve as lemmas that simplify proofs in $T$ by condensing detailed
proof arguments into one rule application.  Each application of a
derivable rule can be expanded locally without affecting the overall
proof structure, which in general is not the case with the application
of an admissible rule.  A derivable rule is also admissible, but in
general the converse does not hold.

The following lemma explains the proof-theoretic status of both types
of rules and clarifies the difference between them.

\begin{lemma} \label{lem:AUX}
  Consider a set of rules $\AUX$.
  \begin{enumerate}[(i)]
  \item 
  Each rule from $\AUX$ is admissible in $T$ iff $T \approx T \cup \AUX$. 

\item Each rule from $\AUX$ is derivable in $T$ iff for all proof
  systems $U$ we have $T \cup U \approx T \cup U \cup \AUX$.

\end{enumerate}
\end{lemma}

So derivability is monotone in the sense that if a rule is derivable
in $T$ then it is also derivable in any proof system extending $T$. In
constrast, simple examples show that admissibility is not monotone.

\begin{proof}
  We only prove $(ii)$.
Below, when choosing a rule $R$ from $\AUX$, we assume that it has the format
\[ 
 \qquad \frac{\varphi_1,\LL,\varphi_k}{\varphi}
\]

\III

\NI
$(\Rightarrow)$ Fix a proof system $U$.  We need to show that in every
proof in $T \cup U \cup \AUX$ each application of a rule from $\AUX$
can be eliminated.  We proceed by induction on the lengths of
proofs. Suppose that in a proof the last rule used was the rule $R$ from $\AUX$.
So by the induction hypothesis
$ \vdash_{T \cup U} \varphi_1, \LL, \vdash_{T \cup U} \varphi_k$.  But $R$
derivable, so 
$\varphi_1, \LL, \varphi_k \vdash_{T \cup U} \varphi$ and hence
$\vdash_{T \cup U} \varphi$ follows.

\III

\NI
$(\Leftarrow)$ Choose a rule $R$ from $\AUX$.
We need to show that
$\varphi_1, \LL, \varphi_k \vdash_T \varphi$, or equivalently
$\vdash_{T \cup U} \varphi$, where $U$ consists of the set of axioms
$\Cu{\varphi_1,\LL,\varphi_k}$.  But the latter claim holds since
$\vdash_{T \cup U \cup \AUX} \varphi$.
\end{proof}

Below we shall need only the left-to-right implications.  In what
follows we shall use the notion of derivability to analyze similar
rules. To make it more visible, given two rules $R_1$ and $R_2$ and a
set of rules $\AUX$, we shall write `$R_1$ can be derived from $R_2$
using $\AUX$' instead of `$R_1$ is derivable in $\AUX \cup R_2$'.

The following lemma will allow us to structure the proof of the main result
in a natural way.

\begin{lemma} \label{lem:equ}
  Suppose that
  \begin{enumerate}[(i)]
  \item $R_1$ can be derived from $R_2$ using $\AUX$,
    
  \item each rule from $\AUX$ is admissible in $T \cup R_2$.    
  \end{enumerate}

  Then $T \cup R_1 \preceq T \cup R_2$.
\end{lemma}

\begin{proof}
  By $(i)$ and Lemma \ref{lem:AUX}$(ii)$
\[
T \cup R_2 \cup \AUX \approx T \cup R_1 \cup R_2 \cup \AUX
\]
and by $(ii)$ and Lemma \ref{lem:AUX}$(i)$
\[
T \cup R_2 \approx T \cup R_2 \cup \AUX.
\]
So 
\[
T \cup R_2 \approx T \cup R_1 \cup R_2 \cup \AUX,
\]
from which the claim follows since
\[
  T \cup R_1 \preceq T \cup R_1 \cup R_2 \cup \AUX.
\]
\end{proof}


\subsection{Auxiliary Rules}
\label{subsec:aux}

The next step consists of identifying the appropriate auxiliary rules. 
In the proof-theoretic analysis of the LOOP II and III rules, we make
use of the following two auxiliary rules, see, e.g., \cite{ABO09}.
\III

\NI
RULE \label{rul:conj} CONJUNCTION
\[ 
 \frac{ \HT{p_1}{S}{q_1}, \HT{p_2}{S}{q_2}     }
      { \HT{p_1 \A p_2}{S}{q_1 \A q_2}         } 
\]
\NI
RULE \label{rul:intro} $\te$-INTRODUCTION 
\[ 
 \frac{ \HT{p}{S}{q}          }
      { \HT{\te z_0:p}{S}{q}    }
\]
\NI
where $z_0$ is a \emph{logical} variable that does not occur  in $\mathit{free}(q)$.
\III

To reason about the LOOP IV rule we adopt the following auxiliary
rules that relate provability in two proof systems.
\III

\NI
RULE PARTIAL CORRECTNESS
\[
 \begin{array}{l}
   \HT{p}{S}{q}\\
  [-\medskipamount]
  \hrulefill                         \\
  \vdash_{\rm{I}} \HT{p}{S}{q}
  \end{array}
\]

\III

\noindent
RULE HYBRID CONJUNCTION
\[
 \begin{array}{l}
  \vdash_{\rm{I}} \HT{p_1}{S}{q_1},   \\
  \HT{p_2}{S}{q_2}                    \\
  [-\medskipamount]
  \hrulefill                          \\
  \HT{p_1 \A p_2}{S}{q_1 \A q_2}
 \end{array}
\]

This rule generalizes the following rule from \cite{ABO09}.
\III

\noindent
RULE DECOMPOSITION
\[
 \begin{array}{l}
  \vdash_{\rm{I}} \HT{p}{S}{q},      \\
  \HT{p}{S}{\T}                      \\
  [-\medskipamount]
  \hrulefill                         \\
  \HT{p}{S}{q}
  \end{array}
\]

\III

% Finally, to reason about the LOOP V rule we use the following
% SUBSTITUTION rule, see, e.g., \cite{ABO09}.

% \III
% \NI
% RULE \label{rul:SUBST} SUBSTITUTION \\
% *** ERO renamed $t$ into $t_0$ to avoid clashes with other $t$ in proofs

% \[ 
%  \frac{ \HT{p}{S}{q}          }
%       { \HT{p[x:=t_0]}{S}{q[x:=t_0]}    }
% \]
% where $x$ is a logical or a program variable
% that does not occur in $S$ and 
% $\mbox{\it var}(t_0)\cap \mbox{\it var}(S)=\emptyset$ ($t_0$ thus may contain logical variables).
% \III

Note that each of these auxiliary rules only \emph{adapts} the assertions without changing the program.
In the context of Hoare's logic such rules are sometimes called \emph{adaptation rules}
\cite{Hoa71,Old83,GR16}.



\subsection{Structure of the Proof}
\label{subsec:structure}

Equivalence Theorem \ref{thm:abo} comprises three equivalences, but
thanks to the transitivity it suffices to prove two of them. We found
it most convenient to deal with the equivalence of the following two
pairs of proof systems: (II, III) and (III, IV). To this end we,
invoke Lemma~\ref{lem:equ} four times, each time establishing its
assumptions for an appropriate pair of LOOP rules $R_1$ and $R_2$ and
a specific set of rules $\AUX$ drawn from the set of auxiliary rules
listed in the previous section, augmented with the CONSEQUENCE rule
(that also is an adaptation rule).

More specifically, we establish the following two lemmas.

\begin{lemma}[Derivability]
  \label{lem:derivability}
The following derivability results hold:

\begin{enumerate} [(i)]
   \item The {\rm LOOP II} rule can be derived from the {\rm LOOP III} rule, using \\
   the  $\exists$-{\rm INTRODUCTION} rule and the {\rm CONSEQUENCE} rule.
   
    \vspace{2mm}

   \item The {\rm LOOP III} rule can be derived from the {\rm LOOP II} rule, using \\
   the {\rm CONJUNCTION} rule and the {\rm CONSEQUENCE} rule.
   
    \vspace{2mm}

   \item The {\rm LOOP IV} rule can be derived from {\rm LOOP III} rule, using \\
   the  {\rm HYBRID CONJUNCTION} rule and the {\rm CONSEQUENCE} rule.
   
    \vspace{2mm}

   \item The {\rm LOOP III} rule can be derived from the {\rm LOOP IV} rule, using \\
   the {\rm PARTIAL CORRECTNESS} rule and the {\rm CONSEQUENCE} rule.
   
%    \vspace{2mm}

   % \item The {\rm LOOP V} rule can be derived from the {\rm LOOP II} rule, using \\
   % the  {\rm SUBSTITUTION} rule and the {\rm CONSEQUENCE} rule.
   
   %  \vspace{2mm}

   % \item The {\rm LOOP II} can be derived from the {\rm LOOP V} rule, using \\
   % the {\rm CONSEQUENCE} rule.
   
\end{enumerate}

\end{lemma}



\begin{lemma}[Admissibility] \label{lem:admissibility}
The following admissibility results hold:


 \begin{enumerate} [$(i)$] 
    
   \item The $\te$-{\rm INTRODUCTION} rule is admissible in the proof system {\rm II}. 
   
    \vspace{2mm}
        
   \item The $\te$-{\rm INTRODUCTION} rule is admissible in the proof system {\rm III}.
   
    \vspace{2mm}
  
   \item The {\rm CONJUNCTION} rule is admissible in the proof system {\rm II}.
   
    \vspace{2mm}

   \item The {\rm HYBRID CONJUNCTION} rule is admissible in the proof system {\rm III}.
%          that is, \\
%          if $\vdash_{\rm{I}} \HT{p_1}{S}{q_1}$ and $\vdash_{\rm{III}} \HT{p_2}{S}{q_2}$, then
%          $\vdash_{\rm{III}} \HT{p_1 \A p_2}{S}{q_1 \A q_2}$.
         
    \vspace{2mm}
         
   \item The {\rm PARTIAL CORRECTNESS} rule is admissible in the proof system {\rm IV}.
   
   %  \vspace{2mm}
   
   % \item The {\rm SUBSTITUTION} rule is admissible in the proof system {\rm II}.

 \end{enumerate}

\end{lemma}

The proofs are given in
Appendix~\ref{sec:proof}.



To close this discussion note the following consequence of the Equivalence Theorem~\ref{thm:abo}
and Lemma \ref{lem:AUX}.

\begin{corollary} \label{corr:admiss}
Each {\rm LOOP} $\alpha$ rule is admissible in the proof system $\beta$, 
        where $\alpha, \beta \in \{\mathrm{II, III, IV}\}$.
\end{corollary}


\begin{proof}
  $T$ be the proof system from Appendix~\ref{sec:I} without the LOOP I
  rule.  By the Equivalence Theorem~\ref{thm:abo} for all LOOP rules
  $\alpha, \beta \in \{\mathrm{II, III, IV}\}$ we have
  $T \cup \alpha \approx T \cup \beta$ and hence
  $T \cup \beta \approx T \cup \{\alpha, \beta\}$.  So by
  Lemma~\ref{lem:AUX}(i), the LOOP $\alpha$ rule is 
  admissible in $T \cup \beta$, which is proof system $\beta$.
% We show this for the case $\alpha$ = II and $\beta$ = III.
% By the Equivalence~\ref{thm:abo}, proof system~II is equivalent to proof system~III,
% so: proof system III $\approx$ proof system~III $\cup$ LOOP II.
% By Lemma~\ref{lem:AUX}(i), LOOP II is admissible in proof system III.
\end{proof}



\subsection{Restrictions on Bound Functions}

A careful reader will notice that in the LOOP II, III and IV rules we
imposed the restriction $var(t) \subseteq \mbox{\it pvar}$ on the
variables of the bound function $t$ that is absent in the publications
that use the LOOP II and III rules.
The reason for this restriction is of technical nature: to be able 
to complete the proofs.
For showing in Lemma~\ref{lem:derivability} that the LOOP~II rule can
be derived from the LOOP~III rule, we need the auxiliary rule
$\exists$-INTRODUCTION. In turn, in Lemma~\ref{lem:admissibility}
we prove that this auxiliary rule is admissible in the proof system III.
This proof is by induction on the structure of {\bf while} programs.
The crucial case is that of a {\bf while} loop. By induction hypothesis,
we can deduce from 
\[
 \vdash_{\rm{III}} \HT{p_{0} \A B \A t = z}{S_0}{t < z}
\]
that
\[
 \vdash_{\rm{III}} \HT{\te z_0: (p_{0} \A B \A t = z)}{S_0}{\te z_0: (t < z)}.
\]
Now, to proceed successfully, we wish to deduce further 
\[
 \vdash_{\rm{III}} \HT{(\te z_0:p_0) \A B \A t = z}{S_0}{t < z}.
\]
This is possible if $z_0 \not\in var(B) \cup var(t)$ and $z_0 \neq z$.
While the last inequality can be assumed w.l.o.g., the first condition
is restrictive.
Since $z_0$ is a logical variable and $var(B) \subseteq \mbox{\it pvar}$,
the imposed restriction $var(t) \subseteq \mbox{\it pvar}$ makes this reasoning
possible.

In a previous version of this paper~\cite{ABO25}, 
without the distinction of logical and program variables,
we required for the bound function $t$ that $var(t) \subseteq var(B) \cup var(S_0)$
holds. For the $\exists$-INTRODUCTION rule it was then required that 
$z_0 \not\in var(B) \cup var(S_0)$.
So again,  $z_0 \not\in var(B) \cup var(t)$.
However, the  condition $var(t) \subseteq var(B) \cup var(S_0)$ is too restrictive 
for some applications. For example, in~\cite{ABO09}, we treat in Section 3.9 the case study
``Partitioning an Array'' and consider a program {\it PART} that corresponds to
the program {\it Partition} in~\cite{FH71}, where Foley and Hoare prove the partial correctness
of Hoare's Quicksort algorithm~\cite{Hoa61b}.
This program has nested while loops. When proving the termination of the inner loops we needed
bound functions that refer to a variable that occurs outside these {\bf while} loops
(see p.110 of~\cite{ABO09}), so the restriction is not satisfied.

In this paper, we distinguish disjoint sets $\mbox{\it lvar}$ and $\mbox{\it pvar}$
of logical variables and program variables, respectively.
A logical variable appears as $z$ to freeze the value of the bound function $t$.
The variables inside programs are all program variables.
A bound function $t$ could contain both logical and program variables.
Our requirement $var(t) \subseteq \mbox{\it pvar}$ is therefore a restriction.
It enables the proof of Lemma~\ref{lem:admissibility}, as explained above.
The restriction is weak enough for performing the termination proof
of the partitioning program {\it PART}.

What about a LOOP rule without any restriction on the bound function $t$\,?
To prove the Equivalence Theorem \ref{thm:abo} in this setting we
reason at the level proofs instead of at the level of proof rules.

Consider a proof of a correctness formula $\phi$ in a proof system
$\alpha \in \{\mathrm{II, III, IV}\}$ in which no distinction is made
between the logical and program variables and in which in the
corresponding LOOP rule no restriction is imposed on the bound
function $t$. First, replace all occurrences of the $z$ variables by
distinct fresh variables that do not occur in the considered program
or in the used bound functions $t$. Clearly, this remains a proof of
$\phi$ in $\alpha$.

Next, let the set of logical variables \emph{lvar} consist of these
fresh variables and let \emph{pvar} be the set of remaining variables.
Then all applications of the corresponding LOOP rule obey the
restriction used in this paper. So by the Equivalence Theorem
\ref{thm:abo} the correctness formula $\phi$ can also be proved in any
other proof system $\beta \in \{\mathrm{II, III, IV}\}$ and hence,
a fortiori in the proof system $\beta$ in which in the corresponding
LOOP rule no restriction on the bound function $t$ is used.


% ****
% 
% By using the following auxiliary SUBSTITUTION rule, see, e.g., \cite{ABO09},
% we can get rid of the remaining restriction.
% 
% \III
% \NI
% RULE \label{rul:SUBST} SUBSTITUTION 
% \[ 
%  \frac{ \HT{p}{S}{q}          }
%       { \HT{p[x:=t_0]}{S}{q[x:=t_0]}    }
% \]
% where $x$ is a logical or a program variable
% that does not occur in $S$ and 
% $\mbox{\it var}(t_0)\cap \mbox{\it var}(S)=\emptyset$ ($t_0$ thus may contain logical variables).
% \III
% 
% The idea is as follows. Suppose a termination proof with LOOP III rule
% requires a bound function $t$ containing logical variables.
% Then apply the substitution rule to replace these variables by fresh program variables.
% This leads to a new bound function $t'$ satisfying the restriction 
% $var(t') \subseteq \mbox{\it pvar}$. Then apply the LOOP III rule as usual.
% At the end, apply the reverse substitution to replace the fresh program variables
% back to the original logical variables.
% What about the status of the SUBSTITUTION rule\,?
% Regrettably, we have not been able to show that it is an admissible in the proof system III
% in the same way that we show it for the other auxiliary rules in Lemma~\ref{lem:admissibility}.
% In an attempted inductive proof, when it comes to the case of the while loop,
% we face the same problems than outlined above that we circumvented by
% restricting the bound function.
% 
% ****




% \section{Soundness and Completeness Matters}

% Assume now that some notion of semantics of programs and assertions is given 
% that includes the concept of a state, execution of a program,
% the notion of a state satisfying an assertion, and the truth of an assertion.
% We write then $\models \HT{p}{S}{q}$ to denote the fact that every execution of $S$ 
% that starts in a state satisfying $p$ terminates in a state satisfying $q$ 
% and say then that $\HT{p}{S}{q}$ is \bfe{true}.
% (Thus we are referring to \bfe{total correctness}.)
% Next, we say then that a proof rule
% \[
%   \frac{\varphi_1,\LL,\varphi_k}{\varphi}
% \]
% is \bfe{sound} if
% $\models \varphi_1,\LL, \models \varphi_k$ implies $\models \varphi$.

% In \cite{ABO09} we proved that the LOOP III rule is sound, while in \cite{Rey98} 
% it was proved that the LOOP II is sound. A natural question arises 
% whether soundness of one of these rules can be directly deduced from
% the soundness of the other rule. 

% \III
% \NI
% ADDED

% This deduction follows from the derivability results stated
% in Lemma~\ref{lem:derivability} % the  main Theorem \ref{thm:main} 
% and the soundness of the
% auxiliary rules.
% For example, suppose now that the LOOP II rule is sound.
% Since the LOOP III rule is derivable from the
% LOOP II rule, its soundness  follows directly from 
%  the soundness of the LOOP II rule and the soundness
%  of the auxiliary rules, in particular of
%  the 
% CONSEQUENCE rule.

% \NI
% END.


% % Further, we can derive for  soundness
% % of the LOOP II rule from the soundness of the LOOP V rule,
% % using the derivability result of clause (vI).
% % However, this approach breaks down for the soundness
% % of the LOOP V rule itself, because Theorem \ref{lem:main} does not
% % provide a corresponding derivability result.
% % The `invasive' notion of admissibility does not allow
% % such a reduction of the soundness of the LOOP V rule
% % to that of the LOOP II rule.
% % Note that the showstopper here is the restriction
% % on the term $t$ in the LOOP rule II,
% % which on the other hand is needed for ther admissibility
% % of the $\exists$-INTRODUCTION rule.

% % This can be accomplished by  modifying the claims of Lemma \ref{lem:1} as follows.

% % \begin{lemma} \label{lem:2}
% % Suppose that $z$ is an integer variable that
% % does not appear in $p,B,t$ or $S$.
% % Then
% %   \begin{enumerate} [(i)]
% %   \item 
% %     If $\models \HT{p \A B} {S}{p}$ and $\models \HT{p \A B \A t=z}{S}{t<z}$, \\ 
% %     then $\models \HT{p \A B \A t=z}{S}{p \A t<z}$.
    
% %   \item
% %     If \ $\models \HT{p \A B \A t=z}{S}{p \A t<z}$, \\ 
% %     then $\models \HT{p \A B} {S}{p}$ and $\models \HT{p \A B \A t=z}{S}{t<z}$.

% %   \end{enumerate}
% % \end{lemma}

% % \begin{proof}
% %   It is a direct consequence of the fact that the proof rules used 
% %   in Lemma~\ref{lem:1} are sound. Thus each time one of these rules
% %   is applied, truth of the correctness formulas is preserved. 
% %   \HB
% % \end{proof}

% % Suppose now that the LOOP II rule is sound. To prove the soundness 
% % of the LOOP III rule assume that its premises are true. 
% % Then by Lemma \ref{lem:2}$(i)$ the premises of the LOOP II rule are true, 
% % so by its soundness the conclusion of both rules is true.
% % The same argument shows that soundness of the LOOP III rule 
% % implies soundness of the LOOP II rule.

% % Finally, for the rules LOOP II--IV we assumed that the bound function $t$ satisfies 
% % the restriction $var(t) \subseteq var(B) \cup var(S)$.
% % This allows for the proof of admissibility of the $\exists$-INTRODUCTION rule 
% % in Theorem~\ref{thm:aux}, which in turn is used via Lemma~\ref{lem:1} 
% % in the proof of Lemma~\ref{lem:main}
% % and thus in the proof of the equivalence result stated in Theorem~\ref{thm:main}.
% % In future we will investigate how to avoid this restriction.
% % We see that the equivalence of the different proof systems for total correctness
% % depends very subtly on the intricate interplay of the loop rules with 
% % the admissibility of standard auxiliary rules of Hoare's logic,
% % which requires further investigation.



% \subsection{Completeness}

% \NI 

% SECTION ON COMPLETENESS ADDED


% Next we discuss  \emph{completeness} of the proof systems II, III, IV, and V.
% Since these proof systems are equivalent, it 
% suffices to prove completeness for one of them.
% We opt here for the proof system V 
% for which
% we consider the standard notion of \emph{relative}
% completeness: For any first-order model $M$ of the assertion language
% we will show that  $M\models\HT{p}{S}{q}$
% implies  $T(M)\vdash_{\rm{V}}\HT{p}{S}{q}$, where
% $T(M)$ denotes the first-order theory of $M$,
% for use in the consequence rule.
% Here  $M\models\HT{p}{S}{q}$ defines the truth of 
% $\HT{p}{S}{q}$ in terms of the model $M$, that is,
% the basic notions of states, executions and satisfaction of the first-order formulas are defined in terms of the model $M$.

% We restrict to models $M$ for which 
% we can express in the assertion language
% the \emph{weakest precondition} $\mbox{\it wp}(S,q)$, for any statement
% $S$ and postcondition $q$.
% This weakest precondition $\mbox{\it wp}(S,q)$ denotes the set
% of initial states for which $S$ terminates in a final state
% satisfying $q$.
% We refer to Section \ref{sec:WP}
% for the standard weakest precondition calculus.
% Further, we restrict to models for which there exists
% % that these models include a well-founded
% % relation $<$ for the interpretation of the corresponding
% % relation symbol in the LOOP V rule, and 
% a parameterized extension $p(x)$
%  of  $\mbox{\it wp}(\WDD{B}{S},q)$, for some   fresh variable $x$
%  not appearing in $\WDD{B}{S}$ or $q$.
% This parameterized extension 
% implies $x\geq 0$ and
% satisfies 
% the following equivalences.

% % \begin{equation}\label{eq:wp1}
% % \mbox{\it wp}(\WDD{B}{S},q) \leftrightarrow
% % \exists x: \mbox{\it wp}(x,\WDD{B}{S},q)   
% % \end{equation}

% \begin{equation}\label{eq:wp1}
% (\exists x: p(x)) \leftrightarrow
% \mbox{\it wp}(\WDD{B}{S},q) 
% \end{equation}


% \begin{equation}\label{eq:wp2}
% p(x)\leftrightarrow
% (((\neg B)\wedge q)\vee (B\wedge \mbox{\it wp}(S, \exists y: p(y) \wedge y<x )))
% \end{equation}

% Instead of representing the bound function by a term $t$
% we thus use in the parameterized extension of the weakest precondition a variable $x$
% to denote an upper-bound of the number of iterations.




% \begin{theorem}\label{thm:comp}
% For any model $M$ for which there exists
% a parameterized extension 
% for every  weakest precondition of a loop statement
% and postcondition, we have that
% $M\models\HT{p}{S}{q}$ implies
% $T(M)\vdash_{\rm{V}}\HT{p}{S}{q}$.
% \end{theorem}

% \begin{proof}
% It suffices to show that
% $T(M)\vdash_{\rm{V}}\HT{\mbox{\it wp}(S,q)}{S}{q}$,
% for any $S$ and postcondition $q$, because $M\models\HT{p}{S}{q}$
% if and only if $p\rightarrow \mbox{\it wp}(S,q)\in T(M)$,
% by definition.
% The proof proceeds by induction on $S$.

% We treat the main case of a loop statement $\WDD{B}{S}$
% (the other cases follow directly from the induction hypotheses using the compositional characterization
% of the weakest precondition described in Section \ref{sec:WP}).
% % For notational convenience, let $p(x)$ abbreviate
% % $\mbox{\it wp}(x,\WDD{B}{S},q)$.
% Let $p(x)$ denote the parameterized extension of
% $\mbox{\it wp}(\WDD{B}{S},q)$.
% By the induction hypothesis
% we have
% \[
% \vdash_{\rm{V}}
% \HT{\mbox{\it wp}(S, \exists y: p(y) \wedge y<x )}{S}{\exists y: p(y) \wedge y<x }
% \]
% Applying the consequence rule, using Equivalence (\ref{eq:wp2}), then gives
% \[
% \vdash_{\rm{V}}
% \HT{p(x)\wedge B}{S}{\exists y: p(y) \wedge y<x }
% \]
% Now we can apply the LOOP V rule and obtain
% \[
% \vdash_{\rm{V}}
% \HT{\exists x:p(x)}{S}{\exists x:p(x)\wedge \neg B}
% \]
% By a final application of the consequence rule, using Equivalence
% (\ref{eq:wp1}) and that $\mbox{\it wp}(\WDD{B}{S},q)\wedge \neg B$ implies
% $q$,
% we obtain the desired result
% \[
% \vdash_{\rm{V}}
% \HT{\mbox{\it wp}(\WDD{B}{S},q)}{S}{q}
% \]

% \III
% \HB
% \end{proof}

% For a deterministic language,  for example the one considered in this paper, 
% the standard weakest precondition in fact amounts to the existence
% of a terminating computation.
% The  parameterized extension of the  standard weakest precondition 
% of a loop statement then additionally expresses
% the number of iterations of 
% the (outer) loop.
% Clearly in this case both the  Equivalences (\ref{eq:wp1})
% and (\ref{eq:wp2}) hold.
% Further, using standard coding techniques, we can show
% the expressibility of the standard weakest precondition and its parameterized
% extension for 
% the standard model of Peano arithmetic.

% \NI
% ADDED A DISCUSSION ON SKOLEM FUNCTIONS.

% Since the proof systems II-V are equivalent,
% completeness of the proof system~V implies completeness
% of the proof systems II-IV.
% However, the proof of this equivalence, i.e., the proof of Theorem
% \ref{thm:main},
% notably between the proof systems II and V, is based on the representation
% of the bound function by  a Skolem function.
% Therefore, to transfer the completeness result for
% the proof system V, we have to extend the notion of expressibility
% by further restricting to
% models which include for every assertion a corresponding
% Skolem function.


% \III
% \NI
% UNBOUNDED NONDETERMINISM: TO BE INCLUDED IN THE SECTION ON EXTENSIONS

% The  Equivalences (\ref{eq:wp1}) 
% and (\ref{eq:wp2}) also hold for  \emph{bounded} non-deterministic  extensions,
% for example, extensions with a non-deterministic choice statement
% $S_1+ S_2$ or Dijkstra's  guarded command language.
% For such languages the  additional parameter then 
% denotes  an \emph{upper-bound} of the number of iterations
% (of the loop statement).

% In programming languages which give rise to
% \emph{unbounded} non-determinism  such an upper-bound may not exist.
% For example
% when the programming language 
% features  a random assignment  which gives
% rise to \emph{countable} non-determinism (see \cite{AptP86}).
% To reason about termination of such programs it suffices
% to consider models which include the initial segment of the \emph{ordinals}
% consisting of all \emph{countable} ordinals.
% For such models the variable $x$ of the LOOP V rule ranges
% over this initial segment, $<$ denotes the membership relation,
% and $0$ denotes the empty set.
% A consequence of this generalization is that we cannot extend the usual coding techniques
% to  express  the standard weakest precondition and its parameterized extension 
% for loop statements, simply because
% the set of all countable ordinals is itself \emph{uncountable}.
% As with the general case of program correctness over 
% abstract data types, one needs more expressive
% assertion languages, e.g., quantification over finite sequences
% (as proposed in \cite{TZ88}).
% Here we resort to an extension of first-order logic
% with a \emph{least fix-point operator}.
% This allows to \emph{define} the weakest precondition $\mbox{\it wp}(S,q)$
% by induction on $S$ and define
% $\mbox{\it wp}(\WDD{B}{S},q)$ by 
% $\exists x: p(x)$, where
% $p(x)$ is defined by
% $$
% \mu R(x,\bar{z}):
% (\neg B)\wedge q)\vee (B\wedge \mbox{\it wp}(S, \exists y:R(y,\bar{z})  \wedge y<x ))
% $$
% Here $\bar{z}$ denote the (free) variables of $\WDD{B}{S}$ and $q$.
% This formula denotes the smallest interpretation of the relation
% $R$ which satisfies 
% $$(\neg B)\wedge q)\vee (B\wedge \mbox{\it wp}(S, \exists y:R(y,\bar{z})  \wedge y<x ))$$
% Such an interpretation exists because this latter formula
% defines a monotonic function on the complete lattice of all
% possible interpretations of $R$, which
% extends any interpretation
% of $R$ by evaluating this latter formula under  this interpretation,
% so that we can apply
% Knaster-Tarski Fixed Point Theorem.
% Further, $\exists x:p(x)$ indeed
% characterizes $\mbox{\it wp}(\WDD{B}{S},q)$
% because 
% we can associate
% with each node of  a terminating computation tree with
% a countable degree an ordinal
% which is bigger than all the ordinals associated with its child nodes.
% Note that the ordinals associated with the child nodes
% form a \emph{set} and therefore there exists an upper-bound
% (and the upper-bound of a countable set of countable ordinals is itself
% a countable ordinal).

% \NI
% END

\section{Representing Proofs}

\subsection{Proof outlines}

Even though the LOOP rules for proving termination are equivalent in
the sense of the Equivalence Theorem~\ref{thm:abo}, they lead to
different proofs of total correctness of \textbf{while} programs.
% This clearly holds for the LOOP V rule which does not involve an
% explicit bound function but uses the invariant for the specification
% of a bound.  But also the other rules lead to different proofs.  
The idea behind the LOOP III rule is to establish that $p$ is a loop
invariant and $t$ is a bound function separately, while in the LOOP II
rule both facts are established simultaneously. So the LOOP III rule
looks more convenient when we want to strengthen a proof of partial
correctness to a proof of total correctness: it suffices to establish
two new premises concerned with the bound function $t$.  In turn, the
LOOP IV rule allows us to split the proof obligations even further, by
identifying the property actually needed to be proved in terms of
total correctness.  This, as already mentioned, allows one to support
modular reasoning.

However, matters change when we want to represent the proofs in the
resulting proof systems  in a convenient form. Given
that these proofs deal with structured programs, their most natural
representation consists of so-called \bfe{proof outlines}, a notion
introduced in \cite{OG76a}.  Informally, it is a proof representation
in the form of a program annotated by the assertions arising from the
appropriate rule applications. Such a representation is possible
thanks to the fact that the proof rules are syntax directed. Proof
outlines were introduced in \cite{OG76a} in order to reason about
correctness of parallel programs, where they served to establish
so-called interference freedom among the proofs of the component
programs. However, they are also very useful as a representation of
correctness proofs of sequential programs and, when some obvious
assertions are deleted, as a program documentation.

Now, given that the LOOP III and IV rules have more than one premise
consisting of a correctness formula, it is difficult to employ
a single proof outline to represent a proof involving any of these rules.
This is not the case with the LOOP II rule. To illustrate this point
recall first that the proof outlines are defined by induction on the program structure.
We only focus on the crucial formation rule concerned with the \textbf{while} statement.
For the proof system I the following formation rule was used in \cite{ABO09}:

\[
\frac{\raisebox{2mm}{$ \HT{p \A B}{S^*}{p}$}}
            {\raisebox{-2mm}{$\HT{{\bf inv:}\ p}
                 {\WDD{B}{\HT{p \A B}{S^*}{p}}}
                 {p \A \neg B}$}}
\]
where $S^*$ is the program $S$ annotated with some assertions.
\III

For the proof system II we introduce the following formation rule:
\III

\NI
\[
\begin{array}{l}
\HT{p \A B \A t=z}{S^{*}}{p \A t<z},                    \\
p \ra t \geq 0                                          \\
[-\medskipamount]
\hrulefill                                              \\
\Cu{{\bf inv:}\ p}
\HT{{\bf bd:}\ t}
   {\WDD{B}{\HT{p \A B \A t = z}{S^*}{p \A t<z}}}
   {p \A \neg B}
\end{array}
\]

\NI
where $S^*$ and $S^{**}$ are annotations of
the program $S$ with some assertions, 
$t$ is an integer expression and $z$ is an integer variable
not occurring in $p,t,B$ or $S^{*}$.\footnote{In \cite{ABO09}, in contrast to \cite{ABO90},
there is a typo and $S^{**}$ is mentioned here instead of $S^{*}$.}
\III


% Finally, for the proof system III we introduce the following formation rule:
% \III

% \NI
% \[
% \begin{array}{l}
% \HT{p \A B \A t=z}{S^{*}}{p \A t<z},                    \\
% p \ra t \geq 0                                          \\
% [-\medskipamount]
% \hrulefill                                              \\
% \Cu{{\bf inv:}\ p}
% \HT{{\bf bd:}\ t}
%    {\WDD{B}{\HT{p \A B \A t = z}{S^*}{p \A t<z}}}
%    {p \A \neg B}
% \end{array}
% \]
% %
% \NI
% where $t$ and $z$ are as above.
% \III

For a moment we defer the discussion of proof outlines for the proof
system~IV.  One can easily prove by induction that each proof outline
for the proof system~I corresponds to a proof in this proof
system. For example, if by the induction hypothesis the proof outline
$\HT{p \A B}{S^*}{p}$ corresponds to a proof of $\HT{p \A B}{S}{p}$ in
I, then the proof outline
\[
  \HT{{\bf inv:}\ p}{\WDD{B}{\HT{p \A B}{S^*}{p}}}{p \A \neg B}
\]
corresponds to a proof of $\HT{p}{\WDD{B}{S}}{p \A \neg B}$ in I,
which is obtained by applying to $\HT{p \A B}{S}{p}$ the LOOP I rule.

However, this property fails to hold for the proof outlines for the proof
system~III, because the second proof outline used in the premise of the
formation rule is dropped. As a consequence, the proof outline for the
\textbf{while} statement does not allow one to reconstruct the proof
of $\HT{p}{\WDD{B}{S}}{p \A \neg B}$ in the proof system III.


\subsection{Proofs using the LOOP II rule}

By contrast, each proof outline for the proof system~II involving the LOOP~II rule
does correspond to a proof in this proof system, because all assertions
used are retained. In other words, from each proof outline for the
proof system~II a proof in this system can be extracted.  

To illustrate this point consider the following program $S_N$ involving nested
loops, suggested to us by Tobias Nipkow (private communication):

\begin{tabbing}
\qquad\qquad
$S_N \equiv$ \= $\WD{i < n}$           \\
\> \qquad $j:=i; $                     \\
\> \qquad $\WD{0 < j}$                 \\
\> \qquad \qquad $j:=j-1$              \\
\> \qquad \OD ;                        \\
\> \qquad $i:=i+1$                     \\
\> \OD                                               
\end{tabbing}
\noindent
where $i, j, n$ are integer variables. 

\III

We would like to prove that it terminates
for all initial states. To this end, we prove the correctness formula
\HT{\T}{S_N}{\T} in the proof system II.
In Fig.~\ref{fig:po}, we show the proof in the form of a proof outline,
instantiating the corresponding formation rule for the \textbf{while} statement in proof system II
with the following loop invariants and bound functions for the outer and the inner loop, respectively:                                                  \\
\[
  \mbox{$p \equiv \T$ and $t \equiv  \mbox{\it max}(n-i,0)$,}
\]
\[
  \mbox{$p \equiv  n-i = z_1 \A z_1>0$ and $t \equiv \mbox{\it max}(j,0)$.}
\]
Note that in a proof outline adjacent assertions stand for implications
according to an application of the CONSEQUENCE rule.
For instance, the assertion in line~3 implies that of line~4.
Assignments are treated by backward substitution according to the ASSIGNMENT axiom.
For instance, the assignment $i:=i+1$ in line~18 is dealt with by 
substituting $i$ by $i+1$ in the assertion in line~19, yielding
the assertion in line~17.


\begin{figure}[ht]

\begin{tabbing}
\ 1\qquad\qquad
 \= \Cu{\textbf{inv} : \T } \Cu{\textbf{bd}: \mbox{\it max}(n-i,0)}    \\
\ 2 \> $\WD{i < n}$                                                  \\
\ 3 \> \qquad \Cu{\T \A i < n \A \mbox{\it max}(n-i,0) = z_1}         \\
\ 4 \> \qquad \Cu{n-i = z_1 \A z_1>0 }                                \\
\ 5 \> \qquad \ \ $j:=i; $                                           \\
\ 6 \> \qquad \Cu{n-i = z_1 \A z_1>0}                                 \\
\ 7 \> \qquad \Cu{\textbf{inv} : n-i = z_1 \A z_1>0} \Cu{\textbf{bd}: \mbox{\it max}(j,0)} \\
\ 8\> \qquad $\WD{0 < j}$                                            \\
\ 9\> \qquad \qquad \Cu{n-i = z_1 \A z_1>0 \A 0<j\A  \mbox{\it max}(j,0)= z_2}  \\
10\> \qquad \qquad \Cu{n-i = z_1 \A z_1>0 \A  j = z_2 \A z_2>0 }      \\
11\> \qquad \qquad \Cu{n-i = z_1 \A z_1>0 \A j-1 < z_2 \A z_2>0 }     \\
12\> \qquad \qquad \ \ $j:=j-1$                                      \\
13\> \qquad \qquad \Cu{n-i = z_1 \A z_1>0 \A j < z_2 \A z_2>0}        \\
14\> \qquad \qquad \Cu{n-i = z_1 \A z_1>0 \A \mbox{\it max}(j,0) < z_2 }        \\
15\> \qquad \OD ;                                                    \\
16\> \qquad \Cu{n-i = z_1 \A z_1>0 \A \neg (0 < j)}                   \\
17\> \qquad \Cu{n-(i+1) < z_1 \A z_1>0 }                              \\
18\> \qquad \ \ $i:=i+1$                                             \\
19\> \qquad \Cu{n-i < z_1 \A z_1>0}                                   \\
20\> \qquad \Cu{\T \A \mbox{\it max}(n-i,0) < z_1 }                   \\
21\> \OD                                                             \\
22\> \Cu{\T \A \neg (i < n)}                                          \\
23\> \Cu{\T}
\end{tabbing}

\caption{Proof outline for \HT{\T}{S_N}{\T} in the proof system II.
The line numbers have been added for reference only.}
\label{fig:po}

\end{figure}



\subsection{Proofs using the LOOP IV rule}

Let us now move on to a discussion of the proofs involving the LOOP IV
rule.  This rule, just like the LOOP III rule, uses more than one
correctness formula as a premise. As a result, it shares with the LOOP
III rule the problem that it is not clear how to faithfully represent
correctness proofs using a single proof outline.
Indeed, each of the first three premises calls for a separate proof
outline.

But it is not difficult to see that this would give rise to largely
overlapping proof outlines.  So, instead, we propose an alternative
approach in which we replace these overlapping proof outlines
referring to the proof system~IV by an interrelated \emph{set} of
proof outlines in the sense of \emph{partial correctness}, so
referring to the proof system~I.

For a given correctness formula $\HT{p}{S}{q}$ to be proved in the
proof system~IV this set is defined as follows.  First, we have a
proof outline $\HT{p}{S^*}{q}$ that employs the first formation rule
given above and thus represents a proof of partial correctness of
$\HT{p}{S}{q}$ in I.

Next, for each occurrence of a loop $\WDD{B}{S_0}$ in $S$, we have a
proof outline $\HT{p_0\wedge B \wedge t=z}{S^*_0}{t<z}$, which
represents a proof of partial correctness of
$\HT{p_0\wedge B \wedge t=z}{S_0}{t<z}$ in I. Here $p_0$ is the
invariant associated with the occurrence of the loop $\WDD{B}{S_0}$ in
the above proof outline $\HT{p}{S^*}{q}$ and $t$ is some bound
function $t$ such that $p_0\to t\geq 0$.  We omit the proof that
existence of such a set of proof outlines in the sense of partial
correctness ensures a proof of the corresponding correctness formula
in the proof system~IV.  Intuitively, the use of the above proof
outlines for \emph{each} loop occurrence in $S$ ensures by structural
induction the third premise of the LOOP~IV rule, so
$\HT{p \A B}{S}{\mathbf{true}}$ in the sense of total correctness.

We illustrate how such a set of proof outlines can be used to
establish the proof of the correctness formula \HT{\T}{S_N}{\T} 
for Nipkow's program $S_N$ in
the proof system~IV.  We skip the trivial proof outline
$\HT{\T}{S^*_N}{\T}$, which corresponds to a proof of partial
correctness of $\HT{\T}{S_N}{\T}$.  Let $S_0$ denote the body of the
outer loop of $S_N$.  Given the bound function
$\mbox{\it max}(n-i,0)$, Fig.~\ref{fig:po2} shows the proof outline
which corresponds to a proof of  partial correctness of
$$\HT{\T\A i < n \A \mbox{\it max}(n-i,0) = z}{S_0}{\mbox{\it max}(n-i,0) < z},$$
assuming that the trivial invariant $\T$ is associated with this loop
in the proof outline $\HT{\T}{S^*_N}{\T}$.  Note that
$\T\to \mbox{\it max}(n-i,0)\geq 0$.
 
 \begin{figure}[ht]
 \begin{tabbing}
 \qquad\qquad                                       
\= \Cu{\T\A i < n \A \mbox{\it max}(n-i,0) = z}  \\
\>  \Cu{n-i = z \A z>0 }                         \\
\>  \ \ $j:=i; $                                \\
\>  \Cu{n-i = z \A z>0}                          \\
\>  \Cu{\textbf{inv} : n-i = z \A z>0}           \\
\>  $\WD{0 < j}$                          \\
\>  \qquad \Cu{n-i = z \A z>0 }                  \\
\>  \qquad \ \ $j:=j-1$                         \\
\> \qquad \Cu{n-i = z \A z>0}                    \\
\>  \OD ;                                       \\
\>  \Cu{n-i = z \A z>0 \A \neg (0 < j)}          \\
\>  \Cu{n-(i+1) < z \A z>0 }                     \\
\>  \ \ $i:=i+1$                                \\
\> \Cu{n-i < z \A z>0}                           \\
\> \Cu{\mbox{\it max}(n-i,0) < z }               \\
[\bigskipamount]
\end{tabbing}
\caption{Proof outline for $\HT{\T\A i < n \A \mbox{\it max}(n-i,0) = z}{S_0}{\mbox{\it max}(n-i,0) < z}$  in the proof system I.}
\label{fig:po2}

\end{figure}

Finally, the LOOP IV rule requires to prove termination of $S_0$.
This boils down to establish the termination of the inner loop.
% Next, we introduce for the inner loop the
To this end, we introduce
the bound function $\mbox{\it max}(j,0)$.
Note that $\T\to \mbox{\it max}(j,0)\geq 0$. 
Fig.~\ref{fig:po3} shows the proof outline
which corresponds to a proof of  partial correctness of
$$\HT{\T\A 0<j\A  \mbox{\it max}(j,0)= z}{j:=j-1}{\mbox{\it max}(j,0) < z}.$$

 \begin{figure}[htbp]
 \begin{tabbing}
 \qquad\qquad    
\=\Cu{\T\A 0<j\A  \mbox{\it max}(j,0)= z}       \\
%\>  \Cu{j = z \A z>0 }                         \\
\>  \Cu{j-1 < z \A z>0 }                        \\
\>  \ \ $j:=j-1$                               \\
\>  \Cu{j < z \A z>0}                           \\
\> \Cu{\mbox{\it max}(j,0) < z }                \\
\end{tabbing}
\caption{Proof outline for $\HT{\T\A 0<j\A  \mbox{\it max}(j,0)= z}{j:=j-1}{\mbox{\it max}(j,0) < z}$  in the proof system I.}
\label{fig:po3}

\end{figure}

% \subsection{Proofs using the LOOP V rule}

% To conclude let us discuss proof outlines in presence of the LOOP V
% rule.  To obtain a proof outline for $\HT{\T}{S_N}{\T}$ in the
% proof system V it suffices to modify the loop invariants in the proof
% outline for the proof system II given in Figure \ref{fig:po}
% % as in the proof of Theorem \ref{thm:main}$(vi)$, so
% by using, respectively,
% \[
%   \mbox{$p(z) \equiv max(n - i, 0) = z$,}
% \]
% \[
%   \mbox{$p(z) \equiv n - i = z_1 \land z_1 > 0 \land max(j, 0) = z$,}
% \]
% and to drop the references to the bound functions.


\section{Practical Applications}

Research on program verification has entered practice most visibly by
the use of assertions as annotations of programs and program
interfaces that remain to be implemented.  This can be seen in the
paradigm of \emph{design by contract} introduced by Bertrand Meyer for
his object-oriented programming language Eiffel~\cite{Mey97}: program
design starts with a specification in terms of assertions, the
contract, against which the program is to be checked either statically
by means of a proof or dynamically at runtime.

This paradigm has been adopted and extended to other programming
languages, in particular to Java.
The \emph{Java Modeling Language} (JML) enriches Java with facilities for 
writing assertions (pre- and postconditions as well as 
class invariants) but also with a concept of abstract state space 
(using so-called model variables)~\cite{JML05}\footnote{See also
  \texttt{https://www.cs.ucf.edu/$\sim$leavens/JML/index.shtml}}.
For assertions, JML uses Java’s Boolean expressions extended by
universal and existential quantifiers. 
% \begin{verbatim}\forall\end{verbatim} and \begin{verbatim}\exists\end{verbatim}.
They are directly written into the Java source code in the form of comments 
starting with the symbols \texttt{//@}
(so that the annotated source code may be processed by both an ordinary Java compiler 
and a specialized JML tool). 

JML is designed to deal with the specification of Java classes,
but here we focus on loops.
JML provides the designated keywords 
\texttt{requires} for specifying the precondition,
\texttt{ensures} for the postcondition,
\texttt{loop\_invariant}, and 
\texttt{loop\_decreases} for the bound function.
To enhance readability, the pre- and postcondition as well as the loop invariant may be split into several
assertions, each one stated after a separate repeated keyword\footnote{For an example, see 
       \texttt{https://www.openjml.org/examples/binary-search.html}}.
An annotated Java program corresponds to a proof outline as discussed here,
but restricted to the essential assertions for each loop (pre- and postcondition,
loop invariant and bound function).

Also for the programming language C, a standardized specification language for C programs, 
called ACSL and inspired by JML, has been designed, see for instance
\cite{BBBCKKMPPS21}\footnote{See also
       \texttt{https://frama-c.com/download/acsl-1.20.pdf}}.

% \texttt{*** For ACSL I found only web pages.}

\medskip

In this paper, we have shown that  the rules LOOP II-IV 
for proving termination of loops are equivalent.
With respect to the number of premises LOOP~II is clearly
the simplest rule for proving termination.
However,  a proof using this rule in general require more
complex assertions because of the accumulation of the
specifications of the bound functions of  inner loops.
This can be seen in the proof outline given in Fig.~\ref{fig:po}, 
in which the assertions inside the inner loop refer to both $z_1$ and $z_2$,
where $z_1$ is the variable freezing the value of the bound function $max(n-i,0)$ of the outer loop
and $z_2$ is the variable freezing the value of the bound function $max(j,0)$ of the inner loop.
The reason is that in this proof outline we need to establish 
that the value of the bound function of the outer loop is not affected 
by the inner loop and that the value of the other bound function decreases.

The LOOP IV rule allows for a separate proof of termination of each
loop, which does not require the specification of bound functions for
the inner loops.  Thus we have a trade off between a single complex
proof and a number of simpler proofs, where complexity is measured by
the size of the assertions.  What works best in practice depends on
the particular program structure.




\section{Proofs from Assumptions}

KRA (5-03-26): NEW SECTION
\III

Admittedly, the equivalence of the proof systems II-IV is a limited
result, since it concerns only the class of \textbf{while}
programs. Here we show how to extend the Equivalence
Theorem~\ref{thm:abo} to a more realistic class of programs.

To this end, consider a situation in which programs can access a library of
procedures each of which is specified by means of a correctness
formula. The access is provided by means of procedure calls.  For
simplicity we assume that procedures do not have global variables and
that Pascal's call by variable mechanism is used.

KRA: SHOULD WE COMMENT ON IT, SO ADD STH LIKE THIS: 

\NI
according to which the formal parameters are initialized at the
beginning of the procedure execution with the values of the actual
parameters, and, at the end of the  procedure's execution are
copied back to the original variables addresses.

\NI ? OR USE SOME OTHER CALL MECHANISM? WHICH ONE?

Let us start with a motivating toy example. Consider the following program
that invokes an unknown (though of course one with the expected meaning)
procedure \emph{swap(x,y)}:

\[
\GCD \equiv \WDD{x \neq q}{\IT{x < y}{swap(x,y)}; x:= x-y}.
\]
KRA: OR SHOULD WE USE PARAMETERS HERE, so $\GCD(x,y)$ ? \\
ERO: I WOULD LEAVE GCD W/O PARAMETERS.

We use in it the statement $\IT{B}{S}$ 
with the corresponding proof rule
\III

\NI
RULE IF
  \[ \frac{ \HT{p \A B}{S}{q}, \ p \A \neg B \ra q}
  { \HT{p}{\IT{B}{S}}{q}}
  \]
\III
  
We want to prove 
\begin{equation}
 \begin{array}{l}  
\HT{x=x_0 \wedge y=y_0}{swap(x,y)}{y=x_0 \wedge x=y_0} \vdash  \\
   \HT{x=x_0 \wedge y=y_0 \wedge x>0 \wedge y>0}{\GCD}{x=y=gcd(x_0, y_0},
 \end{array}
 \label{equ:gcd}
\end{equation}
where $x_0$ and $y_0$ are logical variables and $x$ and $y$ are
program variables.  That is, we want to establish total correctness of
the $\GCD$ program under the stated assumption about the procedure
call \emph{swap(x,y)}.  It is clear how to prove the above claim in
any of the proof systems II-IV: in the proof we can simply use the
assumed property of \emph{swap(x,y)} as an additional axiom. However,
care has to be exercised because some auxiliary rules are needed.

The reason is that we need first to derive a more general property of
\emph{swap(x,y)}, namely the correctness formula
$\HT{p(x,y}{swap(x,y)}{p(y,x)}$, where $p$ is an assertion with free
variables including $x$ and $y$ and which we shall use for a specific
$p$.  To establish it we need the following auxiliary axiom \emph{for
  partial correctness} (see \cite[page 98]{ABO09}):

\III

\NI
AXIOM INVARIANCE
\[ 
\vdash_{\mathrm{I}} \HT{r}{S}{r}
\]
where no variable in $\mathit{free}(r)$ can be changed in $S$.
\III


Using it for the assertion $p(x_0,y_0)$
together with the assumption
\[
  \HT{x=x_0 \wedge y=y_0}{swap(x,y)}{y=x_0 \wedge x=y_0}
\]
we obtain first by the HYBRID CONJUNCTION rule
\[
  \HT{p(x_0,y_0) \wedge x=x_0 \wedge y=y_0}{swap(x,y)}{p(x_0,y_0) \wedge y=x_0 \wedge x=y_0},
\]
from which by the CONSEQUENCE rule and two applications of
the $\te$-INTRODUCTION rule (we assume here that $x_0$ and $y_0$ are logical variables)
we get
\[
  \HT{\te x_0, y_0: (p(x_0,y_0) \wedge x=x_0 \wedge y=y_0)}{swap(x,y)}{p(y,x)},
\]
from which  by the CONSEQUENCE rule $\HT{p(x,y}{swap(x,y)}{p(y,x)}$ follows.

% The following proof outline establishes the desired property of 
% \emph{swap(x,y)} (the second line is justified by two applications of the
% $\te$-INTRODUCTION II rule):
% \III

% \begin{tabbing}
% \qquad\qquad                                       
% \= \Cu{p(x,y)} \\
% \> \Cu{\te x_0, y_0: x=x_0 \wedge y=y_0 \wedge p(x_0,y_0)} \\
% \> \Cu{x=x_0 \wedge y=y_0 \wedge p(x_0,y_0)} \\
% \>  $swap(x,y) $                         \\
% \> \Cu{y=x_0 \wedge x=y_0 \wedge p(x_0,y_0)} \\
% \> \Cu{p(y,x)}
% \end{tabbing}

Then we can exhibit an appropriate proof outline in the proof
system II that at the relevant place uses the just established property of 
\emph{swap(x,y)} for 
\[
  p \equiv gcd(x,y)=gcd(x_0, y_0)  \wedge x>0 \wedge y>0 \wedge z = x+y \wedge x < y.
\]
 \begin{figure}[ht]
 \begin{tabbing}
 \qquad\qquad                                       
\= \Cu{x=x_0 \wedge y=y_0 \wedge x>0 \wedge y>0} \\
\> \Cu{gcd(x,y)=gcd(x_0, y_0)  \wedge x>0 \wedge y>0} \\
\> \Cu{\textbf{inv} : gcd(x,y)=gcd(x_0, y_0)  \wedge x>0 \wedge y>0} \Cu{\textbf{bd}: x+y}    \\
\> $\WD{x \neq y}$                          \\
\>  \qquad \Cu{gcd(x,y)=gcd(x_0, y_0)  \wedge x>0 \wedge y>0 \wedge z = x+y \wedge x \neq y} \\
\>  \qquad $\IF x < y \THEN$                         \\
\> \qquad \qquad \Cu{gcd(x,y)=gcd(x_0, y_0)  \wedge x>0 \wedge y>0 \wedge z = x+y \wedge x < y} \\
\>  \qquad \qquad  $swap(x,y) $                         \\
\> \qquad \qquad \Cu{gcd(y,x)=gcd(x_0, y_0)  \wedge y>0 \wedge x>0 \wedge z = y+x \wedge y < x} \\
\> \qquad \qquad \Cu{gcd(x,y)=gcd(x_0, y_0)  \wedge x>0 \wedge y>0 \wedge z = x+y \wedge x > y} \\
\>  \qquad \FI; \\
\>  \qquad \Cu{gcd(x,y)=gcd(x_0, y_0)  \wedge x>0 \wedge y>0 \wedge z = x+y \wedge x > y} \\
\>  \qquad $x:= x-y$                         \\
\>  \qquad \Cu{gcd(x,y)=gcd(x_0, y_0)  \wedge x>0 \wedge y>0 \wedge z < x+y} \\
\>  \OD                                      \\
\> \Cu{gcd(x,y)=gcd(x_0, y_0)  \wedge x>0 \wedge y>0 \wedge x=y} \\
\> \Cu{x=y = gcd(x,y) \wedge gcd(x,y)=gcd(x_0, y_0)} \\
\> \Cu{x=y=gcd(x_0, y_0)} 
\end{tabbing}
\caption{Proof outline for the conditional correctness formula (\ref{equ:gcd}).}
\label{fig:gcd}

\end{figure}

To extend our study of the LOOP rules to such a setting we
begin by allowing uninterpreted procedure calls, so program statements
of the form $P(\bar{t})$. This leads to a larger class of programs to
which we are referring below. Next, we consider \emph{conditional
  correctness formulas}, which are constructs of the form
\[
\phi_1, \dots, \phi_n \vdash \HT{p}{S}{q},
\]
where $n \geq 0$, each $\phi_i$ is an assertion or a correctness
formula referring to a procedure call and $S$ is a program that may
contain any of the procedure calls appearing in some $\phi_i$.

It is clear how to modify all considered proof rules to deal with the
conditional correctness formulas: it suffices to prefix each
correctness formula by `$\phi_1, \dots, \phi_n \vdash$'.  
% In the case of non-recursive procedures we
% can recover in this setting `unconditional’ correctness formulas by
% adding to each proof system the following proof rules.
% \III

% \NI
% DISCHARGE

%  \[
%  \begin{array}{l}
%   \phi_1 , \dots, \phi_n, \\
%   \phi_1 , \dots, \phi_n \vdash \HT{p}{S}{q} \\
%  [-\medskipamount]
%   \hrulefill                                \\
%   \HT{p}{S}{q}
%  \end{array}
% \]

% SHOULD THE TWO LINES IN THE PREMISES BE SWAPPED TO AVOID A WRONG READING?

% \III

% \NI
% PROCEDURE CALL

% \[
%   \frac{\HT{p}{S}{q}}{\HT{p}{P(\bar{x})}{q}}
% \]
% where $S$ is the body of the procedure $P$ and $\bar{x}$ are the
% formal parameters.
% \III

Additionally, we add the INVARIANCE axiom
\III

\NI
SELECT AXIOM
\[
\phi_1, \dots, \phi_n \vdash \phi_i,
\]
where $i \in \{1, \dots, n\}$.
\III

This yields the proof systems CI-CIV.
Optionally, we could add an appropriate rule that allows us to
derive from assumptions about generic procedure calls adjusted
assumptions about procedure calls with actual parameters.  For details
see the INSTANTIATION rule in \cite[page 159]{ABO09} which lists the
appropriate conditions on formal and actual parameters.

Note that for the \textbf{while} programs the INVARIANCE axiom is admissible in the proof system I, that is the following
observation holds.

\begin{note} \label{note:1}
  For every \textbf{while} program $S$ we have
  $\vdash_{\mathrm{I}} \HT{r}{S}{r}$, where no variable in
  $\mathit{free}(r)$ can be changed in $S$.

\end{note}

However, this axiom cannot be dispensed with for uninterpreted procedure calls, which explains why we
need to retain it.

We would like to extend the Equivalence Theorem~\ref{thm:abo} to the
proof systems CI-CIV, that is to prove their equivalence.  However,
such a result does not hold.  To see this consider the following
example.  For a procedure call $S$ we clearly can prove in CIII
\[
  \begin{array}{l}
    \HT{p \A B}{S}{p}, \ \HT{p \A B \A t=z}{S}{t<z}, \ p \ra t \geq 0 \vdash \\
    \HT{p}{\WDD{B}{S}}{p \A \neg\ B}.
  \end{array}
\]
Indeed, it just suffices to apply the LOOP III rule.
However, this conditional correctness formula cannot be proved in CII
because in this proof system there is no way to prove the premise
\[
  \begin{array}{l}  
    \HT{p \A B}{S}{p}, \ \HT{p \A B \A t=z}{S}{t<z}, \ p \ra t \geq 0 \vdash  \\
    \HT{p \A B \A t=z}{S}{p \A t<z}
  \end{array}    
\]
of the LOOP II rule.

It is worthwhile to note that this premise can be directly proved
using the CONJUNCTION rule.  Now, the Admissibility Lemma
\ref{lem:admissibility}$(iii)$ states that this rule is admissible in
the proof system II, so one might think that such an application of
this rule can be eliminated.  However, in this extended setting this
cannot be done, because we deal with uninterpreted procedure calls.
More specifically, the extended correctness formula
\[
  \HT{p_1}{S}{q_1}, \HT{p_2}{S}{q_2}  \vdash \HT{p_1 \A p_2}{S}{q_1 \A q_2}         
\]
can be proved by applying first the SELECT axiom twice to get 
\[
  \HT{p_1}{S}{q_1}, \HT{p_2}{S}{q_2}  \vdash \HT{p_1}{S}{q_1} 
\]
and
\[
  \HT{p_1}{S}{q_1}, \HT{p_2}{S}{q_2}  \vdash \HT{p_2}{S}{q_2}
\]
and then by applying the CONJUNCTION rule.  If this last rule
application could be eliminated we would establish that the
CONJUNCTION rule is derivable in the proof system II, which is not the
case.

Consequently, only the following counterpart of the Equivalence Theorem~\ref{thm:abo}
holds, where $\AUX$ denotes the set of auxiliary rules introduced in Subection \ref{subsec:aux}.
To prove the Equivalence
Theorem~\ref{thm:abo}, we also use auxiliary rule in the Derivability
Lemma~\ref{lem:derivability} and then show that these auxiliary rules
are admissible in Lemma~\ref{lem:admissibility}.  Why does this
two-step approach not work here as well ?

\begin{theorem} \label{thm:main2}
  The proof systems CII $\cup \AUX$ -- CIV $\cup \AUX$  are equivalent.  
\end{theorem}

\begin{proof}
  We only need to deal with the applications of the LOOP rules. To this end it
  is sufficient to invoke a straightforward modification of the
  Derivability Lemma~\ref{lem:derivability} to the conditional
  correctness formulas obtained by prefixing each correctness formula
  by `$\phi_1, \dots, \phi_n \vdash$'.  By items $(i)$ and $(ii)$ of
  this lemma we get that CII $\cup \AUX$ and CIII $\cup \AUX$ are
  equivalent and by items $(iii)$ and $(iv)$ we get that CIII
  $\cup \AUX$ and CIV $\cup \AUX$ are equivalent. This concludes the
  proof.
\end{proof}

This result underlies, yet again, the importance of auxiliary rules in this comparative study of
the LOOP rules.

%%%%%%%%%%%%%%%%%%%%%%%%%%%%%%%%%%%%%%%%%%%%%%%%%%%%%%%%%%%%%%%%%%%%%%%%%%%%%%
%
\bibliographystyle{ACM-Reference-Format}
% \bibliographystyle{abbrv}
\bibliography{term25}
%
%%%%%%%%%%%%%%%%%%%%%%%%%%%%%%%%%%%%%%%%%%%%%%%%%%%%%%%%%%%%%%%%%%%%%%%%%%%%%%

\appendix

\section{Proof System I}
\label{sec:I}

The proof system I consists of the following axioms and rules:
\VV

\NI
AXIOM \label{rul:skip} SKIP
\[ \HT{p}{skip}{p} \]
%
\NI
AXIOM \label{rul:assi} ASSIGNMENT
\[ \HT{p[u:=t]}{u:=t}{p} \]
%
\NI
RULE \label{rul:comp} COMPOSITION
\[ \frac{ \HT{p}{S_1}{r}, \HT{r}{S_2}{q}                }
        { \HT{p}{S_1;\ S_2}{q}                          }\]
%
\NI
RULE \label{rul:cond} CONDITIONAL
\[ \frac{ \HT{p \A B}{S_1}{q}, \HT{p \A \neg B}{S_2}{q}         }
        { \HT{p}{\ITE{B}{S_1}{S_2}}{q}                          }\]
%
\NI
RULE \label{rul:loop} LOOP I
\[ \frac{ \HT{p \A B}{S}{p}                             }
        { \HT{p}{\WDD{B}{S}}{p \A \neg B}               }\]
%

\NI
RULE \label{rul:cons} CONSEQUENCE
\[ \frac{ p \ra p_1, \HT{p_1}{S}{q_1}, q_1 \ra q        }
        { \HT{p}{S}{q}                                  }\]
%
%\cb

Additionally, given an interpretation $\mathcal{I}$ for the underlying first-order language, 
we use as axioms all assertions that are true in $\mathcal{I}$. 
These assertions are used as premises in the CONSEQUENCE rule. 




% \section{Unrestricted Bound Functions} \label{LOOPIIt}


% % In \cite{ABO09}, to prove termination, we used the LOOP III rule
% % without the stipulation that $var(t) \subseteq var(B) \cup var(S)$.
% % %Call such a rule the LOOP II-$t$ rule.
% % We call the LOOP rules II-IV without this stipulation \emph{unrestricted}.
% % We now prove the following result.
% % % \bigskip


% % \NI
% % RULE \label{rul:loop2} LOOP II-$t$
% % \[
% % \begin{array}{l}
% % \HT{p \A B}{S}{p},                      \\
% % \HT{p \A B \A t=z}{S}{t<z},             \\
% % p \ra t \geq 0                         \\
% % [-\medskipamount]
% % \hrulefill                              \\
% % \HT{p}{\WDD{B}{S}}{p \A \neg\ B}
% % \end{array}
% % \]
% % %
% % \NI
% % where $t$ is an integer expression and $z$ is an integer variable
% % that does not appear in $p,B,t$ or $S$.
% % \III

% \textbf{Proof of Theorem \ref{thm:book}.}

% % \begin{theorem} \label{thm:book}
% % \mbox{} \\
% % The unrestricted  rules LOOP II, III, and IV, are  admissible in the proof system II, III, and IV, respectively.
% % \end{theorem}

% We begin by formulating a crucial lemma.
% Given a list of distinct variables $\bar{x}$ and a list of terms
% $\bar{t}$ of the same length, we denote by $[\bar{x} := \bar{t}]$ the
% \emph{simultaneous substitution} of each variable in $\bar{x}$ by the
% corresponding term in $\bar{t}$. The simultaneous substitution applied
% to a term $t$ or an assertion $p$, written as $t[\bar{x} := \bar{z}]$
% (respectively, $p[\bar{x} := \bar{z}]$), is defined in the usual way
% by properly taking care of the bound variables (see, e.g.,
% \cite[Subsection 2.7]{ABO09}).

% Below, given a list of distinct constants $\bar{c}$ and a list of
% terms $\bar{t}$ of the same length, we denote by
% $[\bar{c} := \bar{t}]$ the \emph{simultaneous replacement} of each
% constant in $\bar{c}$ by the corresponding term in $\bar{t}$, defined
% in the expected way.

% The following lemma formalizes the intuition that in the correctness
% proofs variables that are not used in the analyzed program play the
% same role as fresh constants. The first item concerns a simultaneous
% substitution, while the second one deals with a simultaneous
% replacement.

% \begin{lemma} \label{lem:z}

%   Let \it{PR} be one of the proof systems II, III or IV.  Suppose
%   that $\vdash_{PR} \HT{p}{S}{q}$.

% Assume that $\bar{x}$ and $\bar{c}$ are respectively a list of distinct variables and a list
% of distinct constants of the same length, none of which appears in $S$.
% Then
% \begin{enumerate}[(i)]
% \item $\vdash_{PR} \HT{p[\bar{x} := \bar{c}]}{S}{q[\bar{x} := \bar{c}]}$,

% \item $\vdash_{PR} \HT{p[\bar{c} := \bar{x}]}{S}{q[\bar{c} := \bar{x}]}$.
% \end{enumerate}

% \end{lemma}

% \begin{proof}
%   (Sketch)

% \NI
% The proof proceeds by induction on the structure of $S$ and is straightforward. We illustrate it by presenting
% just one case for $(ii)$, when
% $S \equiv \WDD{B}{S_0}$ and the proof system II. 
% %As in the proof of Theorem \ref{thm:1}$(i)$ [Theorem 2$(i)$ in our paper]
% Let 
% % \[
% %  \vdash_{\rm{II}} \HT{p_{0} \A B}{S_0}{p_{0}},
% % \]
% \[
%  \vdash_{\rm{II}} \HT{p_{0} \A B \A t = z}{S_0}{p_0\wedge  t < z},
% \]
% for some assertion $p_{0}$,  bound function $t$ and
% variable~$z$ such that
% the implications
% $p \ra p_{0}$, $p_{0} \ra t \ge 0$ and  $(p_{0} \A \neg B) \ra q$
% hold.

% Denote $p[\bar{c} := \bar{x}]$ by $p'$ and similarly with
% $p_{0}, B, q$ and $t$.  By the induction hypothesis the above 
% statement holds when we replace each assertion or term by its primed
% version and the primed versions of all three implications hold by
% logic.  Note that by the assumption about the constants in $\bar{c}$
% we have $B' \equiv B$.  Hence by the LOOP II rule
% \[
%   \vdash_{\rm II} \HT{p'_{0}}{\WDD{B}{S}}{p'_{0} \A \neg B},
% \]
% so by the CONSEQUENCE rule
% \[
%   \vdash_{\rm II} \HT{p'}{\WDD{B}{S}}{q'},
% \]
% as desired.
% \HB
% % The desired proof is obtained by performing, respectively, the simultaneous substitution
% % $[\bar{x} := \bar{c}]$ or the simultaneous replacement
% % $[\bar{c} := \bar{x}]$ on all assertions and integer expressions used
% % in the proof of $\HT{p}{S}{q}$ in \emph{PR}.
% \end{proof}

% Below we shall use this lemma twice, first to simultaneously
% substitute specific variables by (fresh) constants and then to
% replace these constants by the original variables.

% We can now establish Theorem \ref{thm:book}. We only treat the case of the unrestricted LOOP II rule that
% was needed to prove item $(iv)$ of Lemma \ref{lem:main}.

% Suppose
% % \[
% %   \vdash_{\rm{II}}  \HT{p \A B}{S}{p},
% % \]
% \[
%   \vdash_{\rm{II}}  \HT{p \A B \A t=z}{S}{p\wedge t<z},
% \]
% and that
% $p \ra t \geq 0$ holds,
% where $t$ is an integer expression such that
% $var(t) \not\subseteq var(B) \cup var(S)$ and $z$ is an integer variable
% that does not appear in $p,B,S, t$ or $q$.

% % \[
% %   \vdash_{\rm{II}}    \HT{p}{\WDD{B}{S}}{p \A \neg\ B},
% % \]
% % where the last step in the proof consists of an application of the
% % LOOP II rule.  We now show that this correctness formula can also be
% % proved in the proof system II when in the LOOP II rule we additionally
% % stipulate that $var(t) \subseteq var(B) \cup var(S)$.

% Let $\bar{x}$ be a list formed by the variables from $var(t)$ that are
% not elements of $var(B) \cup var(S)$ and let $\bar{c}$ be a list of
% distinct constants of the same length as $\bar{x}$, none of which
% appears in $p,B,t$ or $S$. Denote
% $p[\bar{x} := \bar{c}]$ by $p'$ and $t[\bar{x} := \bar{c}]$ by $t'$.

% Then we have by Lemma \ref{lem:z}$(i)$
% % \[
% %  \vdash_{\rm{II}} \HT{p' \A B}{S}{p'},
% % \]
% \[
%  \vdash_{\rm{II}}  \HT{p' \A B \A t'=z}{S}{p'\wedge t'<z},         
% \]
% and
% \[
%   p' \ra t' \geq 0.
% \]

% Note that $var(t') \subseteq var(B) \cup var(S)$. So by an application of the LOOP II rule
% we obtain
% \[
%  \vdash_{\rm{II}}    \HT{p'}{\WDD{B}{S}}{p' \A \neg\ B},
% \]
% from which we get by Lemma \ref{lem:z}$(ii)$
% \[
%  \vdash_{\rm{II}}    \HT{p}{\WDD{B}{S}}{p \A \neg\ B},
% \]
% since $p'[\bar{c} := \bar{x}] \equiv p$ and $B[\bar{c} := \bar{x}] \equiv B$.
% \HB
% \VV

% % Analogous results hold for the same relaxations of the LOOP III and
% % IV rules.
% % Let us explain now the relevance of Theorem \ref{thm:book}.




\section{Proof of the Equivalence Theorem~\ref{thm:abo}}
\label{sec:proof}

As explained in Subsection \ref{subsec:structure}, to prove the
Equivalence Theorem~\ref{thm:abo} (The proof systems II--V are equivalent.),
it suffices to instantiate Lemma~\ref{lem:equ}, 
for each of the four pair of rules
\[
  (R_1, R_2) \in \Cu{(\mathrm{II}, \mathrm{III}),
    (\mathrm{III},\mathrm{II}),
    (\mathrm{III}, \mathrm{IV}),
    (\mathrm{IV}, \mathrm{III})}
\]
and each time for an appropriate set of rules $\AUX$.


\subsection{Derivability of the LOOP Rules}
\label{subsec:proof-derivability}

We provide the proof of the Derivability Lemma~\ref{lem:derivability}, 
repeating each time the statement to be proved.


%\begin{lemma}[Derivability] \label{lem:derivability}

% The following derivability results hold:

% \begin{enumerate}% [(i)]
%    \item The {\rm LOOP II} rule can be derived from the {\rm LOOP III} rule, using \\
%    the  $\exists$-{\rm INTRODUCTION} rule and the {\rm CONSEQUENCE} rule.
   
%     \vspace{2mm}

%    \item The {\rm LOOP III} rule can be derived from the {\rm LOOP II} rule, using \\
%    the {\rm CONJUNCTION} rule and the {\rm CONSEQUENCE} rule.
   
%     \vspace{2mm}

%    \item The {\rm LOOP IV} rule can be derived from {\rm LOOP III} rule, using \\
%    the  {\rm HYBRID CONJUNCTION} rule and the {\rm CONSEQUENCE} rule.
   
%     \vspace{2mm}

%    \item The {\rm LOOP III} rule can be derived from the {\rm LOOP IV} rule, using \\
%    the {\rm PARTIAL CORRECTNESS} rule and the {\rm CONSEQUENCE} rule.
   
% %    \vspace{2mm}

%    % \item The {\rm LOOP V} rule can be derived from the {\rm LOOP II} rule, using \\
%    % the  {\rm SUBSTITUTION} rule and the {\rm CONSEQUENCE} rule.
   
%    %  \vspace{2mm}

%    % \item The {\rm LOOP II} can be derived from the {\rm LOOP V} rule, using \\
%    % the {\rm CONSEQUENCE} rule.
   
% \end{enumerate}


%\end{lemma}

\begin{proof}
The proofs have the same shape given that the last premise of all LOOP rules is the same.
  
\NI
$(i)$ The {\rm LOOP II} rule can be derived from the {\rm LOOP III} rule, using
   the  $\exists$-{\rm INTRODUCTION} rule and the {\rm CONSEQUENCE} rule.
   

It suffices to show how to derive from the premise
\[
  \HT{p \A B \A t=z}{S}{p \A t<z}
\]
of the LOOP II rule the first two premises of the LOOP III rule.  We
first derive by the CONSEQUENCE rule both
\[
 \HT{p \A B \A t=z}{S}{p} \quad \mbox{\rm and} \quad
\HT{p \A B \A t=z}{S}{t<z},
\]
which is the second premise of the LOOP III rule.

Next, by the $\te$-INTRODUCTION rule, we derive 
\[
 \HT{\te z: (p \A B \A t=z)}{S}{p}.
\]
from
$\HT{p \A B \A t=z}{S}{p}$.

By assumption $\mbox{\it var}(t)\subseteq \mbox{\it pvar}$ and $z$ does not appear in $p$. Also
$\mbox{\it var}(B)\subseteq \mbox{\it pvar}$, so
$z$, as a logical variable, does not appear
in $p, B$ or $t$. Hence
\[
  p \A B \to (p \A B \A \te z: t=z) \to  \te z: (p \A B \A t=z).
\]
So by the CONSEQUENCE rule we derive also the first premise
$\HT{p \A B}{S}{p}$ of the LOOP III rule.
\III

\NI
$(ii)$ The {\rm LOOP III} rule can be derived from the {\rm LOOP II} rule, using
   the {\rm CONJUNCTION} rule and the {\rm CONSEQUENCE} rule.
   
It suffices to note that from the first two premises 
\[
  \HT{p \A B} {S}{p}\quad \mbox{\rm and} \quad
   \HT{p \A B \A t=z}{S}{t<z}
\]
of the LOOP III rule we can derive the first premise
\[
  \HT{p \A B \A t=z}{S}{p \A t<z}
\]
of the LOOP II rule by the CONJUNCTION and CONSEQUENCE rules.
\III

\NI
$(iii)$ 
The {\rm LOOP IV} rule can be derived from {\rm LOOP III} rule, using
the  {\rm HYBRID CONJUNCTION} rule and the {\rm CONSEQUENCE} rule.

It suffices to note that from the first three premises 
\[
  \vdash_{\rm{I}}  \HT{p \A B}{S}{p},\quad
  \vdash_{\rm{I}}  \HT{p \A B \A t=z}{S}{t<z},\quad
\mbox{\rm and} \quad
 \HT{p \A B}{S}{\mathbf{true}}
\]
of the LOOP IV rule
we can derive the first two premises
\[
  \HT{p \A B}{S}{p} \quad
  \mbox{\rm and} \quad \HT{p \A B \A t=z}{S}{t<z}
\]
of the LOOP III rule by two applications of the HYBRID CONJUNCTION and
CONSEQUENCE rules.
\III

\NI
$(iv)$
The {\rm LOOP III} rule can be derived from the {\rm LOOP IV} rule, using 
   the {\rm PARTIAL CORRECTNESS} rule and the {\rm CONSEQUENCE} rule.

From the first two premises
\[
  \HT{p \A B}{S}{p}
\quad
\mbox{\rm and} \quad
   \HT{p \A B \A t=z}{S}{t<z},
\]
of the LOOP III rule
we obtain the first two premises
$$\vdash_{\rm{I}}  \HT{p \A B}{S}{p}
\quad
\mbox{\rm and} \quad \vdash_{\rm{I}}  \HT{p \A B \A t=z}{S}{t<z}$$
of the LOOP IV rule
by applying the PARTIAL CORRECTNESS rule.
The third premise
\[
  \HT{p \A B}{S}{\mathbf{true}}
\]
of the LOOP IV rule is obtained by the CONSEQUENCE rule.
\end{proof}
% \NI
% $(v)$ This is the only involved case in the proof. It assumes that in the underlying first-order
% language for all assertions appropriate \emph{Skolem functions} exist. More precisely, we postulate
% that for each assertion $p(x)$ with the free variables $x, \bar{u}$ there exists 
% a \emph{Skolem function}   $f(\bar{u})$ such that the implication
% \[
%   (\exists x:p(x))\rightarrow p(f(\bar{u}))
% \]
% holds. Let
% \[
%   p_{\min}(x) \equiv p(x) \land \fa y: (p(y) \ra x \leq y).
% \]

% Then given the premises
% \[
%  \HT{p(x)\wedge B}{S}{\exists y: p(y)\wedge y< x}
% \]
% and $p(x)\rightarrow x\geq 0$ of the LOOP V rule, where $x, \bar{u}$
% are the free variables of $p$, we want to apply the LOOP II rule with
% $\exists x:p(x)$ as the loop invariant and $t \equiv f(\bar{u})$ as
% the bound function, where $f(\bar{u})$ is the Skolem function for
% $p_{\min}(x)$.
% % such that
% % $(\exists x:p(x))\rightarrow p(f(\bar{u}))$
% % and $p(x)\rightarrow f(\bar{u})\leq x $. 
% Intuitively, $t$ denotes the minimal witness for $p(x)$, which exists
% because by the second premise of the rule the implication
% $p(x) \ra \te y: p_{\min}(y)$ holds.

% However, some of the variables in $\bar{u}$ may be logical, so this
% definition of $t$ does not satisfy the condition
% $\mbox{\it var}(t)\subseteq\mbox{\it pvar}$ imposed on $t$.  Therefore
% we first apply the SUBSTITUTION rule to obtain
% \[
%  \HT{p'(x)\wedge B}{S}{\exists y: p'(y)\wedge y< x}
% \]
% where $p'(x)$ is the result of substituting every logical variable in
% $p(x)$ (excluding $x$) by a fresh \emph{program} variable.  Let
% $x, \bar{u}$ denote the list of free variables of $p'(x)$ and let
% $f(\bar{u})$ be the Skolem function for $p'_{\min}(x)$.

% Because of the second premise of the LOOP V rule and the definitions
% of $p'_{\min}(x)$ and of the Skolem function, the following two strings of implications hold:
% \[
% ((\exists x:p'(x)) \wedge f(\bar{u})=x) \ra p'_{\min}(x) \ra p'(x)
% \]
% and
% \[
%   \begin{array}{l}
%   (\exists y: p'(y) \wedge y < x) \ra (\exists y: p'_{\min}(y)
%     \wedge y < x) \ra \\
%     (p'_{\min}(f(\bar{u})) \wedge f(\bar{u}) < x)
%   \ra ((\te x: p'(x)) \wedge f(\bar{u}) < x).
%   \end{array}
% \]
% An application of the CONSEQUENCE rule then gives
% \[
%  \HT{(\exists x:p'(x))\wedge B\wedge f(\bar{u})=x}{S}{(\exists x:p'(x))\wedge f(\bar{u})<x}.
% \]
% Further, by the second premise of the LOOP V rule we have
% $p'(x)\rightarrow x\geq 0$ which implies
% $p'_{\min}(x)\rightarrow x\geq 0$.  Hence
% \[
%   (\exists x:p'(x)) \ra (\exists x:p'_{\min}(x)) \ra f(\bar{u})\geq 0.
% \]
% We can now apply 
% the LOOP II rule with $\exists x:p'(x)$ as the loop invariant and
% $t \equiv f(\bar{u})$ as the bound function to obtain 
% \[
%   \HT{\exists x: p'(x)}{\WDD{B}{S}}{(\exists x:p'(x))  \A \neg B}.
% \]
% The desired conclusion
% \[
%   \HT{\exists x: p(x)}{\WDD{B}{S}}{(\exists x:p(x))  \A \neg B}
% \]
% now follows by an application of the SUBSTITUTION rule.
% \III

% \NI
% KRA: (13-02-26) SLIGHTLY REWRITTEN

% \NI
% $(vi)$ Given the premises
% \[
%  \HT{p\wedge B\wedge t=z}{S}{p\wedge t<z}
% \]
% and $p\rightarrow t\geq 0$ of the LOOP II rule let
% $p'(z)$ denote the assertion $p\wedge t=z$.
% For any fresh variable $y$ the implication
% \[
%  (p\wedge t<z)\rightarrow \exists y: (p\wedge t=y \wedge y<z)
% \]
% holds. By the CONSEQUENCE rule we thus obtain
% \[
%  \HT{p'(z)\wedge B}{S}{\exists y: (p'(y)\wedge y< z)}.
% \]
% Further, $p\rightarrow t\geq 0$ implies
% $p'(z)\rightarrow z\geq 0$.
% So applying the LOOP V rule we obtain 
% \[
%  \HT{\exists z:p'(z)}{\WDD{B}{S}}{(\exists z:p'(z)) \wedge  \neg B}.
% \]
% Since  the equivalence $(\exists z:p'(z))\leftrightarrow p$ holds, applying
% the CONSEQUENCE rule
% yields the desired conclusion
% \[
%   \HT{p}{\WDD{B}{S}}{p \wedge \neg B}.
% \]
% \HB




\subsection{Admissibility of the Auxiliary Rules}
\label{subsec:proof-admissibility}

We show the admissibility of each auxiliary rule in the proof system
with the version of the LOOP rule that used this auxiliary rule 
according to Lemma~\ref{lem:derivability}. 
Additionally, we prove statement $(i)$ of Lemma~\ref{lem:admissibility} 
because it is needed in the proof of statement $(iii)$.
Since the CONSEQUENCE rule is included in all these proof systems,
its admissibility is trivial.


% \begin{lemma}[Admissibility] \label{lem:admissibility}
% The following admissibility results hold:


%  \begin{enumerate} [$(i)$] 
    
%    \item The $\te$-{\rm INTRODUCTION} rule is admissible in the proof system {\rm II}. 
   
%     \vspace{2mm}
        
%    \item The $\te$-{\rm INTRODUCTION} rule is admissible in the proof system {\rm III}.
   
%     \vspace{2mm}
  
%    \item The {\rm CONJUNCTION} rule is admissible in the proof system {\rm II}.
   
%     \vspace{2mm}

%    \item The {\rm HYBRID CONJUNCTION} rule is admissible in the proof system {\rm III}.
% %          that is, \\
% %          if $\vdash_{\rm{I}} \HT{p_1}{S}{q_1}$ and $\vdash_{\rm{III}} \HT{p_2}{S}{q_2}$, then
% %          $\vdash_{\rm{III}} \HT{p_1 \A p_2}{S}{q_1 \A q_2}$.
         
%     \vspace{2mm}
         
%    \item The {\rm PARTIAL CORRECTNESS} rule is admissible in the proof system {\rm IV}.
   
%    %  \vspace{2mm}
   
%    % \item The {\rm SUBSTITUTION} rule is admissible in the proof system {\rm II}.

%  \end{enumerate}

% \end{lemma}

We first establish two lemmas that provide additional
information about the proofs in the considered proof systems.

\begin{lemma} \label{lemma:1}
Let $\mathit{PR}$ be one of the proof systems I, II, III, or IV.
If $\mathit{PR} \vdash \varphi$, there exists a proof of $\varphi$ in $\mathit{PR}$ 
with exactly one final application of the CONSEQUENCE rule.
\end{lemma}

\begin{proof}
Since implication $\ra$ is reflexive, we can always add to a given proof in $\mathit{PR}$
one final application of the CONSEQUENCE rule.
Since implication is transitive, successive applications of the CONSEQUENCE rule in a proof in $\mathit{PR}$
can be condensed into one application.
\end{proof} 

\begin{lemma} \label{lem:useful}
Let $\mathit{PR}$ be one of the proof systems I, II, III, or IV.
  \begin{enumerate} [(i)] 
  \item  Suppose that  $\vdash_{PR} \HT{p}{S_1; \ S_2}{q}$.
    Then for some assertion $r$
    \[
 \vdash_{PR} \HT{p}{S_1}{r} \mbox{ and }
 \vdash_{PR} \HT{r}{S_2}{q}.
\]

\item Suppose that  $\vdash_{PR} \HT{p}{\ITE{B}{S_1}{S_2}}{q}$.
    Then 
    \[
 \vdash_{PR} \HT{p \A B}{S_1}{q} \mbox{ and }
 \vdash_{PR} \HT{p \A \neg B}{S_2}{q}.
\]
 \end{enumerate}    
\end{lemma}
\begin{proof}
  \mbox{} \\
  $(i)$ By Lemma~\ref{lemma:1}, the considered correctness formula was proved using
  the COMPOSITION rule followed by a single application of the CONSEQUENCE rule.
  So for some assertions $p_1, r, q_1$, we have
    \[
 \vdash_{PR} \HT{p_1}{S_1}{r} \mbox{ and }
 \vdash_{PR} \HT{r}{S_2}{q_1}
\]
and the implications $p \to p_1$ and $q_1 \to q$ hold.
We now get the claim by the CONSEQUENCE rule.
\III

\NI
$(ii)$ The argument is analogous as in $(ii)$. 
\end{proof}

We now turn to the proof of the Admissibility Lemma~\ref{lem:admissibility}, 
repeating each time the statement to be proved.

\begin{proof}
 \mbox{} \\
 
%%%%%%%%%%%%%%%%%%%%%%%%%%%%%%%%%%%%%%%%%%%%%%%%%%%%%%%%%%%%%%%%
\NI
$(i)$  The $\te$-{\rm INTRODUCTION} rule is admissible in the proof system {\rm II}.

We proceed by induction on the structure of $S$
and consider a proof of \HT{p}{S}{q} in the proof system \rm{II}, where 
the last two steps involve an axiom or a rule of the proof system \rm{II} for 
the top-level operator of the program $S$, 
followed by one final application of the CONSEQUENCE rule
according to Lemma~\ref{lemma:1}.
By the assumptions on $z_0$, we have $z_0 \not\in var(S) \cup \mathit{free}(q)$.

\III

\NI
$\bullet$ \emph{Case} $S \equiv u:=t$.
Thus suppose $\vdash_{\rm{II}} \HT{p}{S}{q}$.
Then for some assertion $q_1$,
\[
 \vdash_{\rm{II}} \HT{q_1[u:=t]}{S}{q_1}
\]
by the ASSIGNMENT axiom,
and the implications 
$p \ra q_1[u:=t]$ and $q_1 \ra q$ hold.
Since $z_0 \in lvar$, we have $z_0 \not \equiv u$. 
By the assignment axiom, also
\[
 \HT{(\te z_0:q_1)[u:=t]}{S}{\te z_0:q_1}.
\]
Note that the implications 
$\te z_0:p \ra \te z_0: (q_1[u:=t])$ and $\te z_0: q_1 \ra \te z_0: q$ hold.
Since $z_0 \not \equiv u$ and $z_0 \not\in \mathit{free}(q)$, 
also the implications 
\[
 \te z_0: (q_1[u:=t]) \ra (\te z_0:q_1)[u:=t] \ \mbox{ and }\  (\te z_0:q) \ra q
\] 
hold.
So the CONSEQUENCE rule yields
$\vdash_{\rm{II}}  \HT{\te z_0: p}{S}{q}$, as desired.

\III

\NI
$\bullet$ \emph{Case} $S \equiv S_1 ;\ S_2$.
Thus suppose $\vdash_{\rm{II}} \HT{p}{S}{q}$.
By Lemma \ref{lem:useful}$(i)$
for some assertion $r$, 
\[
 \vdash_{\rm{II}} \HT{p}{S_1}{r}  \mbox{ and }
 \vdash_{\rm{II}} \HT{r}{S_2}{q}.
\]
Since $r \ra \te z_0:r$, the CONSEQUENCE rule yields
$\vdash_{\rm{II}}  \HT{p}{S_1}{\te z_0: r}$.
Since $z_0 \not \in \mathit{free}(\te z_0: r)$, the induction hypothesis yields
$\vdash_{\rm{II}}  \HT{\te z_0: p}{S_1}{\te z_0: r}$.
Since $z_0 \not\in \mathit{free}(q)$, 
by the induction hypothesis, also $\vdash_{\rm{II}}  \HT{\te z_0: r}{S_2}{q}$.
Thus by the COMPOSITION  rule, 
$\vdash_{\rm{II}}  \HT{\te z_0: p}{S}{q}$, as desired.


\III

\NI
$\bullet$ \emph{Case} $S \equiv \ITE{B}{S_1}{S_2}$.
Thus suppose $\vdash_{\rm{II}} \HT{p}{S}{q}$.
Then by Lemma \ref{lem:useful}$(ii)$,
\[
 \vdash_{\rm{II}} \HT{p \A B}{S_1}{q} \mbox{ and }
 \vdash_{\rm{II}} \HT{p \A \neg B}{S_2}{q}.
\]
By the induction hypothesis,
\[
 \vdash_{\rm{II}} \HT{\te z_0: (p \A B)}{S_1}{q}  \mbox{ and }
 \vdash_{\rm{II}} \HT{\te z_0: (p \A \neg B)}{S_2}{q}.
\]
Since $z_0 \not\in var(B)$, the CONSEQUENCE rule yields
\[
 \vdash_{\rm{II}} \HT{(\te z_0: p) \A B}{S_1}{q}  \mbox{ and }
 \vdash_{\rm{II}} \HT{(\te z_0: p) \A \neg B}{S_2}{q}.
\]
By the CONDITIONAL rule,
$\vdash_{\rm{II}} \HT{\te z_0: p}{S}{q}$,
as desired.

\III

\NI
$\bullet$ \emph{Case} $S \equiv \WDD{B}{S_0}$.
Thus suppose $\vdash_{\rm{II}} \HT{p}{S}{q}$.
By the assumption about $t$, we have
$z_0 \not \in var(t)$ and without loss of generality we can assume $z_0 \neq z$.
Then for some assertion $p_{0}$ and an appropriate bound function $t$ and
variable~$z$,
\[
 \vdash_{\rm{II}} \HT{p_{0} \A B \A t = z}{S_0}{p_0\wedge t < z},
\]
and the implications
$p \ra p_{0}$, $p_{0} \ra t \ge 0$,  $(p_{0} \A \neg B) \ra q$
hold.

Since $(p_0\wedge t<z) \ra ((\te z_0: p_0)\wedge t<z) $, the CONSEQUENCE rule yields
\[
 \vdash_{\rm{II}} \HT{p_{0} \A B \A t=z}{S_0}{(\te z_0: p_{0})\wedge t<z }.
\]
By the induction hypothesis, the assumption about $z_0$,
and the CONSEQUENCE rule 
\[
 \vdash_{\rm{II}} \HT{(\te z_0:p_0) \A B \A t = z}{S_0}{(\te z_0:p_0)\A t < z}.
\]
So the LOOP II rule yields 
\[
 \vdash_{\rm{II}} \HT{\te z_0:p_0}{S}{(\te z_0:p_0) \A \neg B}.
\]
Further, since $(p_{0} \A \neg B) \ra q$ holds and
$z_0 \not\in (\mathit{free(q)} \cup var(B))$, also the implication
$((\te z_0:p_0) \A \neg B) \ra q$ holds.  So a final application of the
CONSEQUENCE rule yields $\vdash_{\rm{II}} \HT{\te z_0:p}{S}{q}$, as
desired.

\III

%%%%%%%%%%%%%%%%%%%%%%%%%%%%%%%%%%%%%%%%%%%%%%%%%%%%%%%%%%%%%%%%
\NI
$(ii)$ The $\te$-{\rm INTRODUCTION} rule is admissible in the proof system {\rm III}.

Again, we proceed by induction on the structure of $S$,
but now consider a proof of \HT{p}{S}{q} in the proof system \rm{III}, where 
the last two steps involve the axiom or rule of the proof system \rm{III} 
for the top-level operator of the program $S$, 
followed by one final application of the CONSEQUENCE rule
according to Lemma~\ref{lemma:1}.

Except for the \textbf{while} statement, all cases are analogous, 
with $\vdash_{\rm{II}}$ replaced by $\vdash_{\rm{III}}$. 
The case of the \textbf{while} statement differs from $(i)$ only 
in that one now considers two correctness formulas in the premise 
of the LOOP III rule instead of one. 
As above, let $S \equiv \WDD{B}{S_0}$ and
suppose $\vdash_{\rm{III}} \HT{p}{S}{q}$.
By the assumption about $t$, we have
$z_0 \not \in var(t)$ and without loss of generality we can assume $z_0 \neq z$.
Then for some assertion $p_{0}$ and an appropriate bound function $t$ and
variable~$z$,
\[
 \vdash_{\rm{III}} \HT{p_{0} \A B}{S_0}{p_{0}},
\]
\[
 \vdash_{\rm{III}} \HT{p_{0} \A B \A t = z}{S_0}{t < z},
\]
and the implications
$p \ra p_{0}$, $p_{0} \ra t \ge 0$,  $(p_{0} \A \neg B) \ra q$
hold.

Since $p_0 \ra \te z_0: p_0$, the CONSEQUENCE rule yields
\[
 \vdash_{\rm{III}} \HT{p_{0} \A B}{S_0}{\te z_0: p_{0}}.
\]
By the induction hypothesis, the assumption about $z_0$,
and the CONSEQUENCE rule both
\[
 \vdash_{\rm{III}} \HT{(\te z_0:p_0) \A B}{S_0}{\te z_0: p_{0}}
\]
and
\[
 \vdash_{\rm{III}} \HT{(\te z_0:p_0) \A B \A t = z}{S_0}{t < z}.
\]
So the LOOP III rule yields 
\[
 \vdash_{\rm{III}} \HT{\te z_0:p_0}{S}{(\te z_0:p_0) \A \neg B}.
\]
Further, since $(p_{0} \A \neg B) \ra q$ holds and
$z_0 \not\in (\mathit{free(q)} \cup var(B))$, also the implication
$((\te z_0:p_0) \A \neg B) \ra q$ holds.  So a final application of the
CONSEQUENCE rule yields $\vdash_{\rm{III}} \HT{\te z_0:p}{S}{q}$, as
desired.

\III

%%%%%%%%%%%%%%%%%%%%%%%%%%%%%%%%%%%%%%%%%%%%%%%%%%%%%%%%%%%%%%%%
\NI
$(iii)$ The {\rm CONJUNCTION} rule is admissible in the proof system {\rm II}.

We proceed by induction on the structure of $S$
and for $i=1,2$ consider proofs of \HT{p_i}{S}{q_i} in the proof system \rm{II}, where 
the last two steps involve the axiom or rule of the proof system \rm{II} 
for the top-level operator of the program $S$, 
followed by one final application of the CONSEQUENCE rule
according to Lemma~\ref{lemma:1}.

\III

\NI
$\bullet$ \emph{Case} $S \equiv u:=t$.
Thus suppose $\vdash_{\rm{II}} \HT{p_1}{S}{q_1}$ and
$\vdash_{\rm{II}} \HT{p_2}{S}{q_2}$.
Then for some assertions $q_{01}$ and $q_{02}$,
by the ASSIGNMENT axiom, both
\[
 \vdash_{\rm{II}} \HT{q_{01}[u:=t]}{S}{q_{01}} \ \mbox{ and }\
 \vdash_{\rm{II}} \HT{q_{02}[u:=t]}{S}{q_{02}},
\]
and the implications 
$p_1 \ra q_{01}[u:=t]$, $q_{01} \ra q_1$ and
$p_2 \ra q_{02}[u:=t]$, $q_{02} \ra q_2$ hold.
The ASSIGNMENT axiom also yields
\[
 \HT{q_{01}[u:=t] \A q_{02}[u:=t]}{S}{q_{01} \A q_{02}}.
\]
By the implications 
$(p_1 \A p_2) \ra (q_{01}[u:=t] \A q_{02}[u:=t])$ and 
$(q_{01} \A q_{02}) \ra$ $(q_1 \A q_2)$,
the CONSEQUENCE rule yields
\[
  \vdash_{\rm{II}} \HT{p_1 \A p_2}{S}{q_1 \A q_2},
\]
as desired.

\III

\NI
$\bullet$ \emph{Case} $S \equiv S_1 ;\ S_2$.
Thus suppose $\vdash_{\rm{II}} \HT{p_1}{S}{q_1}$ and
$\vdash_{\rm{II}} \HT{p_2}{S}{q_2}$.
Then by Lemma \ref{lem:useful}$(i)$ for some assertions $r_1$ and $r_2$
\[
 \vdash_{\rm{II}} \HT{p_{1}}{S_1}{r_1} \ \mbox{ and }\ 
 \vdash_{\rm{II}} \HT{r_1}{S_2}{q_{1}},
\]
\[
 \vdash_{\rm{II}} \HT{p_{2}}{S_1}{r_2} \ \mbox{ and }\
 \vdash_{\rm{II}} \HT{r_2}{S_2}{q_{2}}.
\]
By the induction hypothesis,
\[
 \vdash_{\rm{II}} \HT{p_1 \A p_2}{S_1}{r_1 \A r_2} \ \mbox{ and }\
 \vdash_{\rm{II}} \HT{r_1 \A r_2}{S_1}{q_1 \A q_2},
\]
so by the COMPOSITION rule,
\[
  \vdash_{\rm{II}} \HT{p_1 \A p_2}{S_1;\ S_2}{q_1 \A q_2},
\]
as desired.

\III

\NI
$\bullet$ \emph{Case} $S \equiv \ITE{B}{S_1}{S_2}$.
Thus suppose $\vdash_{\rm{II}} \HT{p_1}{S}{q_1}$ and
$\vdash_{\rm{II}} \HT{p_2}{S}{q_2}$.
Then by Lemma \ref{lem:useful}$(ii)$,
\[
 \vdash_{\rm{II}} \HT{p_{1}\A B}{S_1}{q_{1}} \ \mbox{ and }\
 \vdash_{\rm{II}} \HT{p_{1}\A \neg B}{S_2}{q_{1}},
\]
\[
 \vdash_{\rm{II}} \HT{p_{2} \A B}{S_1}{q_{2}} \ \mbox{ and }\
 \vdash_{\rm{II}} \HT{p_{2} \A \neg B}{S_2}{q_{2}}.
\]
By the induction hypothesis,
\[
 \vdash_{\rm{II}} \HT{p_{1}\A p_{2} \A B}{S_1}{q_{1} \A q_{2}}
\]
and
\[
 \vdash_{\rm{II}} \HT{p_{1}\A p_{2} \A \neg B}{S_2}{q_{1} \A q_{2}}.
\]
So by the CONDITIONAL rule,
\[
  \vdash_{\rm{II}} \HT{p_1 \A p_2}{\ITE{B}{S_1}{S_2}}{q_1 \A q_2},
\]
as desired.

\III

\NI
$\bullet$ \emph{Case} $S \equiv \WDD{B}{S_0}$.
Suppose $\vdash_{\rm{II}} \HT{p_1}{S}{q_1}$ and
$\vdash_{\rm{II}} \HT{p_2}{S}{q_2}$.
Then for some assertions $p_{01},p_{02}$, appropriate bound functions $t_1, t_2$ and
variables $z_1, z_2$,
\[
 \vdash_{\rm{II}} \HT{p_{01} \A B \A t_1 = z_1}{S_0}{p_{01} \A t_1 < z_1},
\]
\[
 \vdash_{\rm{II}} \HT{p_{02} \A B \A t_2 = z_2}{S_0}{p_{02} \A t_2 < z_2},
\]
the implications
$p_{01} \ra t_1 \ge 0, \, p_1 \ra p_{01}, \, (p_{01} \A \neg B) \ra q_1$
and
$p_{02} \ra t_2 \ge 0,$ $p_2 \ra p_{02}, \, (p_{02} \A \neg B) \ra q_2$
hold.

Without loss of generality we can assume that $z_2 \not \in \{z_1\} \cup \mathit{free}(p_{01})$.
By the induction hypothesis,
\[
 \vdash_{\rm{II}} \HT{p_{01} \A p_{02} \A B \A t_1 = z_1 \A t_2 = z_2}{S_0}
                     {p_{01} \A p_{02} \A t_1 < z_1 \A t_2 < z_2}.
\]

It suffices to consider one bound function, say $t_1$.
Formally, we show this as follows.
By the CONSEQUENCE rule,
\[
 \vdash_{\rm{II}} \HT{p_{01} \A p_{02} \A B \A t_1 = z_1 \A t_2 = z_2}{S_0}
                     {p_{01} \A p_{02} \A t_1 < z_1}.
\]
Now, an application of the $\te$-INTRODUCTION rule, which is admissible
in the proof system II according to part $(i)$ of this lemma, followed by an application
of the CONSEQUENCE rule yields
\[
 \vdash_{\rm{II}} \HT{p_{01} \A p_{02} \A B \A t_1 = z_1 \A \te z_2: t_2 = z_2}{S_0}
                     {p_{01} \A p_{02} \A t_1 < z_1}.
\]
A further application of the CONSEQUENCE rule yields
\[
 \vdash_{\rm{II}} \HT{p_{01} \A p_{02} \A B \A t_1 = z_1}{S_0}
                     {p_{01} \A p_{02} \A t_1 < z_1}.
\]
Then by the LOOP II rule,
\[
  \vdash_{\rm{II}} \HT{p_{01} \A p_{02}}{\WDD{B}{S_0}}{p_{01} \A p_{02} \A \neg B}.
\]
The implications above yield
$(p_1 \A p_2) \ra (p_{01} \A p_{02})$ and
$(p_{01} \A p_{02} \A \neg B) \ra$ $(q_1 \A q_2)$.
Thus by the CONSEQUENCE rule,
\[
 \vdash_{\rm{II}} \HT{p_1 \A p_2}{\WDD{B}{S_0}}{q_1 \A q_2},
\]
as desired.
\III

%%%%%%%%%%%%%%%%%%%%%%%%%%%%%%%%%%%%%%%%%%%%%%%%%%%%%%%%%%%%%%%%
\NI
$(iv)$ The {\rm HYBRID CONJUNCTION} rule is admissible in the proof system {\rm III},
that is,
if $\vdash_{\rm{I}} \HT{p_1}{S}{q_1}$ and $\vdash_{\rm{III}} \HT{p_2}{S}{q_2}$, then
$\vdash_{\rm{III}} \HT{p_1 \A p_2}{S}{q_1 \A q_2}$.

The proof is analogous to the proof of $(iii)$. 
The only case that is somewhat different is the one
concerned with the \textbf{while} statement. 
So we only deal with
\III

\NI
$\bullet$ \emph{Case} $S \equiv \WDD{B}{S_0}$.
Suppose $\vdash_{\rm{I}} \HT{p_1}{S}{q_1}$ and
$\vdash_{\rm{III}} \HT{p_2}{S}{q_2}$.
Then for some assertions $p_{01},p_{02}$ and an appropriate bound function $t$ and
variable~$z$,
\[
 \vdash_{\rm{I}} \HT{p_{01} \A B}{S_0}{p_{01}},
\]
\[
 \vdash_{\rm{III}} \HT{p_{02} \A B}{S_0}{p_{02}},
\]
\[
 \vdash_{\rm{III}} \HT{p_{02} \A B \A t = z}{S_0}{t < z},
\]
and the implications
\[
  p_1 \ra p_{01}, \, (p_{01} \A \neg B) \ra q_1, \,
  p_{02} \ra t \ge 0, \, p_2 \ra p_{02}, \, (p_{02} \A \neg B) \ra q_2
\]
hold. By the induction hypothesis,
\[
 \vdash_{\rm{III}} \HT{p_{01} \A p_{02} \A B}{S_0}
                     {p_{01} \A p_{02}},
\]
and by the induction hypothesis combined with the CONSEQUENCE rule,
\[
 \vdash_{\rm{III}} \HT{p_{01} \A p_{02} \A B \A t = z}{S_0}
                     {t < z}.
\]
Since $p_{01} \A p_{02} \ra t \ge 0$, the LOOP III rule yields
\[
  \vdash_{\rm{III}} \HT{p_{01} \A p_{02}}{\WDD{B}{S_0}}{p_{01} \A p_{02} \A \neg B}.
\]
The implications above yield
$(p_1 \A p_2) \ra (p_{01} \A p_{02})$ and
$(p_{01} \A p_{02} \A \neg B) \ra$ $(q_1 \A q_2)$.
Thus by the CONSEQUENCE rule,
\[
 \vdash_{\rm{III}} \HT{p_1 \A p_2}{\WDD{B}{S_0}}{q_1 \A q_2},
\]
as desired.

\III

%%%%%%%%%%%%%%%%%%%%%%%%%%%%%%%%%%%%%%%%%%%%%%%%%%%%%%%%%%%%%%%%
\NI
$(v)$ The {\rm PARTIAL CORRECTNESS} rule is admissible in the proof system {\rm IV}.

Suppose $\vdash_{\rm{IV}} \HT{p}{S}{q}$. 
Consider a corresponding  proof of \HT{p}{S}{q} in the proof system \rm{IV}.
Each application of the LOOP IV rule requires the establishment of four premises,
of which the first one is exactly the premise an application of the LOOP I rule
depends on. If we eliminate in the given proof of \HT{p}{S}{q} in system \rm{IV}
all parts establishing the last three premises in each application of the
LOOP IV rule, we obtain a proof of \HT{p}{S}{q} in the proof system I 
using only the LOOP I rule. Thus $\vdash_{\rm{I}} \HT{p}{S}{q}$. 
\end{proof}


\subsection{Instantiating Lemma~\ref{lem:equ}}
\label{subsec:instantiate}

Finally, we explain how to obtain the Equivalence Theorem~\ref{thm:abo} by instantiating Lemma~\ref{lem:equ}
using the Derivability Lemma~\ref{lem:derivability} and the Admissibility Lemma~\ref{lem:admissibility}.
We show one illustrative case: proof system II $\approx$ proof system III.

Let $T$ be the proof system from Appendix~\ref{sec:I} without the LOOP I rule.
Then, using Lemma~\ref{lem:derivability}(i) and Lemma~\ref{lem:admissibility}(ii), 
instantiating Lemma~\ref{lem:equ} yields
\[
 T \cup \mathrm{LOOP\ II}\ \preceq\ T \cup \mathrm{LOOP\ III}.
\]
Also, using Lemma~\ref{lem:derivability}(ii) and Lemma~\ref{lem:admissibility}(iii), 
instantiating Lemma~\ref{lem:equ} yields
\[
 T \cup \mathrm{LOOP\ III}\ \preceq T \cup \mathrm{LOOP\ II}.
\]
So altogether $T \cup \mathrm{LOOP\ II}\ \approx\ T \cup \mathrm{LOOP\ III}$,
which means proof system II $\approx$ proof system III.


\section{Admissibility of the LOOP $B$ versions}
\label{sec:B}

*** ERO: NEW VERSION -- FIRST EQUIVALENCE ***

We establish here the following result that clarifies the
proof-theoretic status of the $B$ versions of the LOOP rules.
Let us mark the proof systems using these versions with a $B$ in the end.

\begin{theorem}\label{thm:B}
  Each proof system $\alpha B$ is equivalent to the proof system
  $\alpha$, where $\alpha \in \{{\rm II, III, IV}\}$.
\end{theorem}

\begin{corollary}
  Each {\em LOOP $\alpha B$} rule is admissible in the proof system
  $\alpha$, where $\alpha \in \{{\rm II, III, IV}\}$.
\end{corollary}

We prove Theorem~\ref{thm:B} in detail for the case $\alpha=$ III.
First we show:

\begin{lemma}\label{lem:IIIB}
  The {\rm LOOP III} B rule can be derived from the {\rm LOOP III} using
  the {\rm HYBRID CONJUNCTION} rule,
  the {\rm $\exists$-INTRODUCTION} rule, and the {\rm CONSEQUENCE} rule.
\end{lemma}

\begin{proof}
%   We prove the claim for the LOOP III $B$ rule. The proofs for the
%   other two proof rules are analogous, though they require proving
%   that the HYBRID CONJUNCTION and $\te$-INTRODUCTION rules are
%   admissible in the proof systems II and IV.
%
Suppose that for some loop invariant $p$ and bound function $t$
\begin{equation} \label{eq:inv}
\HT{p \A B}{S}{p},
\end{equation}
\begin{equation} \label{eq:bd}
\HT{p \A B \A t=z}{S}{t<z}
\end{equation}
and
$p \A B \ra t \geq 0$.
By the INVARIANCE axiom in the proof system I (see Note \ref{note:1})
we have
\[
\vdash_{\rm{I}}  \HT{z' = z+1 \A 0 < z'}{S}{z' = z+1 \A 0 < z'},
\]
where $z'$ is a fresh logical variable.
Applying the HYBRID CONJUNCTION rule to this correctness formula and (\ref{eq:bd}),
we obtain
\begin{equation}\label{eq:bd-hc}
 \HT{z' = z+1 \A 0 < z' \A p \A B \A t=z}{S}{z' = z+1 \A 0 < z' \A t<z}.
\end{equation}

Next, note that for the bound function $t' \equiv \ITE{B}{t+1}{0}$
the following implications hold:
\[
  z' = z+1 \A 0 < z' \A p \A B \A t'=z' \ra
    z' = z+1 \A 0 < z' \A p \A B \A t=z
\]
and
\[
  z' = z+1 \A 0 < z' \A t<z \ra t' < z',
\]
so applying the CONSEQUENCE rule to (\ref{eq:bd-hc}) we obtain
\[
  \HT{z' = z+1 \A 0 < z' \A p \A B \A t'=z'}{S}{t'<z'}.
\]

Applying next the $\te$-INTRODUCTION rule,  we get
\begin{equation}\label{eq:bd-te}
  \HT{\te z: (z' = z+1 \A 0 < z' \A p \A B \A t'=z')}{S}{t'<z'}.
\end{equation}

Note now the following sequence of implications:
\[
  \begin{array}{l}
    p \A B \A t'=z' \ra \\
    0 \leq t \A p \A B \A t'=z' \ra \\
    0 \leq t \A t' = t+1 \A p \A B \A t'=z' \ra \\
    0 < z' \A p \A B \A t'=z' \ra \\
    (\te z: z' = z+1) \A 0 < z' \A p \A B \A t'=z'  \ra \\
    \te z: (z' = z+1 \A 0 < z' \A p \A B \A t'=z'). 
  \end{array}
\]

An application of the CONSEQUENCE rule to (\ref{eq:bd-te}) then yields
\begin{equation}\label{eq:bd-new}
  \HT{p \A B \A t'=z'}{S}{t'<z'},
\end{equation}
Further, we have $p \A B \ra t \geq 0 \A t' = t+1$ and 
$p \A \neg B \ra t' = 0$, so 
$p \ra t' \geq 0$ holds.
This allows us to apply the LOOP III rule to (\ref{eq:inv}), (\ref{eq:bd-new}),
and the last implication to establish
\[
 \HT{p}{\WDD{B}{S}}{p \wedge \neg B}.
\]
\end{proof}

By Lemma~\ref{lem:admissibility}, the auxiliary rules needed in Lemma~\ref{lem:IIIB}
are admissible in proof system III. The converse of Lemma~\ref{lem:IIIB}
is immediate.

\begin{lemma}\label{lem:III}
  The {\rm LOOP III} rule can be derived from the {\rm LOOP III B}. 
\end{lemma}

\begin{proof}
Suppose that for some loop invariant $p$ and bound function $t$
\[
\HT{p \A B}{S}{p}, \
\HT{p \A B \A t=z}{S}{t<z}, \mbox{  and }  p \ra t \geq 0.
\]
Since the last condition implies $p \A B \ra t \geq 0$,
an application of the LOOP III B rule yields
\[
 \HT{p}{\WDD{B}{S}}{p \wedge \neg B}.
\]
\end{proof}

Finally, by instantiating Lemma~\ref{lem:equ} with the 
Lemmas~\ref{lem:IIIB}, \ref{lem:III}, and \ref{lem:admissibility}, 
we obtain that the proof system III B is equivalent to the proof system III.
For the cases II and IV, the equivalence proofs require
lemmas corresponding to Lemma~\ref{lem:IIIB} and Lemma~\ref{lem:III},
which can be proven analogously,
though it has to be shown along the lines of Lemma~\ref{lem:admissibility}
that the HYBRID CONJUNCTION rule is admissible in the proof systems II and IV,
and that the $\te$-INTRODUCTION rule is admissible in the proof system IV.



%%%%%%%%%%%%%%%%%%%%%%%%%%%%%%%%%%%%%%%%%%%%%%%%%%%%%%%%%%%%

\III

*** PREVIOUS VERSION: ONLY ADMISSIBILTY ***

\begin{theorem}
  Each LOOP $\alpha B$ rule is admissible in the proof system
  $\alpha$, where $\alpha \in \{II, III, IV\}$.
\end{theorem}

\begin{proof} of the Theorem
  We prove the claim for the LOOP III$B$ rule. The proofs for the
  other two proof rules are analogous, though they require proving
  that the HYBRID CONJUNCTION and $\te$-INTRODUCTION rules are
  admissible in the proof systems II and IV.

  Suppose that for some loop invariant $p$ and bound function $t$
\[
\vdash_{\rm{III}}  \HT{p \A B}{S}{p},
\]
\[
\vdash_{\rm{III}}  \HT{p \A B \A t=z}{S}{t<z}
\]
and
$p \A B \ra t \geq 0$.

By the admissibility of the INVARIANCE axiom in the proof system I (see Note \ref{note:1})
we have
\[
\vdash_{\rm{I}}  \HT{z' = z+1 \A 0 < z'}{S}{z' = z+1 \A 0 < z'},
\]
where $z'$ is a fresh logical variable.

Applying the HYBRID CONJUNCTION rule that by Lemma~\ref{lem:admissibility}$(iv)$
is admissible in the proof system III
we obtain
\[
  \vdash_{\rm{III}}  \HT{z' = z+1 \A 0 < z' \A p \A B \A t=z}{S}{z' = z+1 \A 0 < z' \A t<z}.
\]

Next, note that for the bound function $t' \equiv \ITE{B}{t+1}{0}$
the following implications hold:
\[
  z' = z+1 \A 0 < z' \A p \A B \A t'=z' \ra
    z' = z+1 \A 0 < z' \A p \A B \A t=z
\]
and
\[
  z' = z+1 \A 0 < z' \A t<z \ra t' < z',
\]
so using the CONSEQUENCE rule we obtain
\[
  \vdash_{\rm{III}}  \HT{z' = z+1 \A 0 < z' \A p \A B \A t'=z'}{S}{t'<z'}.
\]

Applying next the $\te$-INTRODUCTION rule that by Lemma~\ref{lem:admissibility}$(ii)$
is admissible in the proof system III we get
\[
  \vdash_{\rm{III}}  \HT{\te z: (z' = z+1 \A 0 < z' \A p \A B \A t'=z')}{S}{t'<z'}.
\]

Note now the following sequence of implications:
\[
  \begin{array}{l}
    p \A B \A t'=z' \ra \\
    0 \leq t \A p \A B \A t'=z' \ra \\
    0 \leq t \A t' = t+1 \A p \A B \A t'=z' \ra \\
    0 < z' \A p \A B \A t'=z' \ra \\
    (\te z: z' = z+1) \A 0 < z' \A p \A B \A t'=z'  \ra \\
    \te z: (z' = z+1 \A 0 < z' \A p \A B \A t'=z'). 
  \end{array}
\]

An application of the CONSEQUENCE rule then yields
\[
\vdash_{\rm{III}}  \HT{p \A B \A t'=z'}{S}{t'<z'},
\]
Further, we have $p \A B \ra t \geq 0 \A t' = t+1$ and 
$p \A \neg B \ra t' = 0$, so 
$p \ra t' \geq 0$ holds.
This allows us to apply the LOOP III rule to establish
\[
\vdash_{\rm{III}}  \HT{p}{\WDD{B}{S}}{p \wedge \neg B}.
\]
\end{proof}

%%%%%%%%%%%%%%%%%%%%%%%%%%%%%%%%%%%%%%%%%%%%%%%%%%%%%%%%%%%%%%%%%%%%%%%%%%%%%%%%
%
\end{document}
%
%%%%%%%%%%%%%%%%%%%%%%%%%%%%%%%%%%%%%%%%%%%%%%%%%%%%%%%%%%%%%%%%%%%%%%%%%%%%%%%%



